\chapter{电子与光子的量子交换}\label{chap:quantum-electron-photon-exchange}

经典外场只给电子一个确定的驱动。若光场本身具有数涨落，电子能量交换还会留下光子记录；联合态与条件测量因而必须从同一个幺正作用求出。

\section{量子交换：多模一般动力学方程与单模闭式}

经典场模型把光的作用写成给定复振幅，适合平均场主导的相干光。若要比较不同光子数分布，或求电子与光之间的纠缠，就需要保留光场的量子态。我们从\chapref{chap:general-interaction}的多模联合哈密顿量出发，先写出一般交换方程，再在单模及低反冲条件下求解析传播子。

多模连续谱层级~\eqnrefs{eq:universal_multiboson_amplitude_hierarchy} 给出量子交换的母动力学。先把有限组已谱隔离的环境模式写成多模联合格；只有当其它模式的有限时间跃迁幅度均低于目标精度时，才保留一个孤立模式。随后独立检查电子逐边失谐与耦合的平移对称性，满足条件后才能引入双边电子平移算符 $\hat b$。光场的相干态大占据极限发生在另一 Hilbert 空间方向，不替代前两项谱选择条件。


量子交换的母模型不应从单模开始。先假设环境谱已经由上一章的一般 Green/正常模分析选出 $M$ 个需要显式保留的正交玻色模式，但电子仍保留完整正能连续谱。第 $\nu$ 个模式的中心频率为 $\omega_\nu$，湮灭算符为 $\hat a_\nu$。在相互作用绘景中，最一般的线性单量子交换仍是连续核
\begin{equation}
\boxed{
\hat H_{{\rm Q},I}^{(M)}(t)
=
\hbar
\sum_{\nu=1}^{M}
\sum_{ss'}
\int\dd^3p\,\dd^3p'\,
\left[
\mathcal G_{\nu;s's}(\mathbf p',\mathbf p;t)
|\mathbf p',s'\rangle\langle\mathbf p,s|
\,\hat a_\nu
+\mathrm{H.c.}
\right].}
\label{eq:multimode_qpinem_continuous_kernel_dp101}
\end{equation}
$\mathcal G_{\nu;s's}$ 由 QED 电流--模式矩阵元或相应材料量子化给出；这一层还没有电子“格点”，也没有为每个光学模式预先指定唯一动量步长。

很多纳米光子学几何中，第 $\nu$ 个模式只在一个窄的纵向 Fourier 区间内与准直电子强耦合。若实际电子工作窗内的主导转移为 $\hbar\kappa_\nu$，并且自旋翻转及其它动量转移的累计传播修正很小，可以在通道子空间上写
\[
\hat{\mathcal G}_{\nu,I}(t)
=
\ii\upsilon_\nu(t)\hat b_\nu^\dagger
+\hat R_{\nu,\kappa}(t),
\qquad
\hat b_\nu^\dagger|p\rangle
=
|p+\hbar\kappa_\nu\rangle .
\]
这里 $\hat R_{\nu,\kappa}$ 收集非主导 Fourier 分量、逐边矩阵元变化和被舍弃自旋通道。若 $\Pi_{\rm in}$ 投影到实际入射工作窗，以只保留主导平移项的传播子记为 $\hat U_{\rm tr}$，则一个直接的有限时间充分条件是
\begin{equation}
\boxed{
\epsilon_{\kappa}^{(M)}
\equiv
\frac1\hbar
\int_{t_i}^{t_f}\dd t\,
\left\|
\hbar\sum_\nu
\left[
\hat R_{\nu,\kappa}(t)\hat a_\nu
+\hat R_{\nu,\kappa}^\dagger(t)\hat a_\nu^\dagger
\right]
\hat U_{\rm tr}(t,t_i)\Pi_{\rm in}
\right\|
\ll1 .}
\label{eq:multimode_momentum_step_projection_error_dp101}
\end{equation}
只有通过这一类通道投影后，才得到常用的多模平移 Hamiltonian
\begin{equation}
\boxed{
 \hat H_{\rm Q}^{(M)}(t)
 =\ii\hbar\sum_{\nu=1}^{M}
 \left[
 \upsilon_\nu(t)\hat b_\nu^\dagger\hat a_\nu
 -\upsilon_\nu^*(t)\hat b_\nu\hat a_\nu^\dagger
 \right].}
\label{eq:multimode_qpinem_hamiltonian_dp47}
\end{equation}
因此 $\upsilon_\nu$ 由连续耦合核在选定电子通道上的投影唯一确定。

重复交换以后，整数向量
\[
\boldsymbol\ell=(\ell_1,\ldots,\ell_M)\in\mathbb Z^M
\]
映射到真实电子动量
\begin{equation*}
\mathbf p_{\boldsymbol\ell}
=
\mathbf p_0
+\hbar\sum_\nu\ell_\nu\boldsymbol\kappa_\nu .
\end{equation*}
把 $|\boldsymbol\ell\rangle_e$ 作为正交的多维格点标签，还要求整数映射在实际占据窗口 $\mathcal W_\ell$ 内保持一一对应。若存在不同整数向量满足
\[
\sum_\nu
(\ell_\nu-\ell_\nu')\boldsymbol\kappa_\nu=0,
\]
它们代表同一个电子动量末态，必须先在振幅层面相干相加，而不能作为不同格点分别计数。对可唯一标记的通道，真实电子能量为
\begin{equation}
 \mathcal E_{\boldsymbol\ell}
 =E\!\left(
 \mathbf p_0+\hbar\sum_\nu
 \ell_\nu\boldsymbol\kappa_\nu
 \right),
\end{equation}
并相对于理想多频能量梯定义
\begin{equation}
\boxed{
 \delta_{\boldsymbol\ell}
 =
 \frac{
 \mathcal E_{\boldsymbol\ell}-E_0
 -\sum_\nu\ell_\nu\hbar\omega_\nu
 }{\hbar}.}
\label{eq:multimode_recoil_detuning_dp47}
\end{equation}

对初始 Fock 扇区
$\mathbf n=(n_1,\ldots,n_M)$，取守恒基底
\[
 |\boldsymbol\ell;\mathbf n\rangle
 =
 |\boldsymbol\ell\rangle_e
 \otimes
 \bigotimes_\nu
 |n_\nu-\ell_\nu\rangle_\nu ,
\]
其中每个光子占据数都必须非负。把联合态展开为
\[
|\Psi_{\mathbf n}(t)\rangle
=
\sum_{\boldsymbol\ell}
c_{\boldsymbol\ell}^{(\mathbf n)}(t)
|\boldsymbol\ell;\mathbf n\rangle ,
\]
逐项使用玻色梯矩阵元后得到
\begin{equation}
\boxed{
\begin{aligned}
 \ii\dot c_{\boldsymbol\ell}^{(\mathbf n)}
 =&\,\delta_{\boldsymbol\ell}
 c_{\boldsymbol\ell}^{(\mathbf n)}\\
 &+\ii\sum_\nu \upsilon_\nu(t)
 \sqrt{n_\nu-\ell_\nu+1}\,
 c_{\boldsymbol\ell-\mathbf e_\nu}^{(\mathbf n)}\\
 &-\ii\sum_\nu \upsilon_\nu^*(t)
 \sqrt{n_\nu-\ell_\nu}\,
 c_{\boldsymbol\ell+\mathbf e_\nu}^{(\mathbf n)} .
\end{aligned}}
\label{eq:multimode_qpinem_lattice_dp47}
\end{equation}
这才是有限组显式模式下的多模 QPINEM 数值一般动力学方程：真实电子色散、多个 Fock 下界、有限脉冲和模式间共享的电子反冲全部保留。单模模型只需令 $M=1$，但若进一步要得到双边位移闭式，还必须额外满足逐边失谐与逐边耦合近似均匀。

\begin{proposition}[多模固定交换记录的联合振幅格]
设有限模主导转移均由电子平移算符表示，并保留各模式的初始光子数记录 $\mathbf n$。基态 $|\mathbf p_{\boldsymbol\ell}\rangle_e\otimes|\mathbf n-\boldsymbol\ell\rangle_{\rm ph}$ 的振幅满足\eqnrefs{eq:multimode_qpinem_lattice_dp47}，连续核的平移约化误差为\eqnrefs{eq:multimode_momentum_step_projection_error_dp101}。即使不同指标给相同电子动量，不同光子记录仍可令联合基态正交；只有末联合记录相同时才相干合并路径。
\end{proposition}
\begin{proof}
先把第 $\nu$ 个连续电子耦合算符记为
\[
\hat{\mathcal G}_{\nu,I}(t)
=
\sum_{ss'}
\int\dd^3p\,\dd^3p'\,
\mathcal G_{\nu;s's}(\mathbf p',\mathbf p;t)
|\mathbf p',s'\rangle\langle\mathbf p,s|.
\]
在实际窄带工作窗内，假定主导矩阵元集中在一个自旋保持、近固定动量转移通道
\[
(\mathbf p,s_0)
\longrightarrow
(\mathbf p+\hbar\boldsymbol\kappa_\nu,s_0).
\]
定义参考平移算符
\[
\hat b_\nu^\dagger
=
\int\dd^3p\,
|\mathbf p+\hbar\boldsymbol\kappa_\nu,s_0\rangle
\langle\mathbf p,s_0|
\]
以及投影后的共同速率 $\upsilon_\nu(t)$。于是精确耦合算符在该工作窗上分解成
\[
\hat{\mathcal G}_{\nu,I}(t)
=
\ii\upsilon_\nu(t)\hat b_\nu^\dagger
+
\hat R_{\nu,\kappa}(t).
\]
把所有余项收集为
\[
\hat R_\kappa(t)
=
\hbar
\sum_\nu
\left[
\hat R_{\nu,\kappa}(t)\hat a_\nu
+
\hat R_{\nu,\kappa}^\dagger(t)\hat a_\nu^\dagger
\right].
\]
令 $\hat H_{\rm tr}$ 只保留主导平移项，并由它生成 $\hat U_{\rm tr}$。完整传播子 $\hat U$ 与 $\hat U_{\rm tr}$ 的 Duhamel 恒等式为
\[
(\hat U-\hat U_{\rm tr})\Pi_{\rm in}
=
-\frac{\ii}{\hbar}
\int_{t_i}^{t_f}\dd t\,
\hat U(t_f,t)
\hat R_\kappa(t)
\hat U_{\rm tr}(t,t_i)\Pi_{\rm in}.
\]
完整传播子幺正，因此
\[
\|(\hat U-\hat U_{\rm tr})\Pi_{\rm in}\|
\le
\frac1\hbar
\int_{t_i}^{t_f}\dd t\,
\|
\hat R_\kappa(t)
\hat U_{\rm tr}(t,t_i)\Pi_{\rm in}
\|,
\]
这正是正文\eqnrefs{eq:multimode_momentum_step_projection_error_dp101}。只有该量足够小时，连续核才受控退化为
\[
\hat H_{\rm Q}^{(M)}
=
\ii\hbar\sum_\nu
\left[
\upsilon_\nu\hat b_\nu^\dagger\hat a_\nu
-\upsilon_\nu^*\hat b_\nu\hat a_\nu^\dagger
\right].
\]

设整数向量 $\boldsymbol\ell$ 对应电子动量（不要求这个电子映射单射）
\[
\mathbf p_{\boldsymbol\ell}
=
\mathbf p_0
+\hbar\sum_\nu\ell_\nu\boldsymbol\kappa_\nu,
\]
并定义
\[
|\boldsymbol\ell;\mathbf n\rangle
=
|\mathbf p_{\boldsymbol\ell}\rangle_e
\otimes
\bigotimes_\nu
|n_\nu-\ell_\nu\rangle_\nu .
\]
对第 $\nu$ 个吸收顶角，
\begin{align*}
\hat b_\nu^\dagger\hat a_\nu
|\boldsymbol\ell;\mathbf n\rangle
&=
\hat b_\nu^\dagger
|\mathbf p_{\boldsymbol\ell}\rangle
\otimes
\sqrt{n_\nu-\ell_\nu}\,
|n_\nu-\ell_\nu-1\rangle_\nu
\otimes_{\mu\neq\nu}
|n_\mu-\ell_\mu\rangle_\mu\\
&=
\sqrt{n_\nu-\ell_\nu}\,
|\boldsymbol\ell+\mathbf e_\nu;\mathbf n\rangle .
\end{align*}
同理，
\[
\hat b_\nu\hat a_\nu^\dagger
|\boldsymbol\ell;\mathbf n\rangle
=
\sqrt{n_\nu-\ell_\nu+1}\,
|\boldsymbol\ell-\mathbf e_\nu;\mathbf n\rangle .
\]

在旋去理想多频能量
$E_0+\sum_\nu\ell_\nu\hbar\omega_\nu$
的表象中，真实电子色散留下对角项
\[
\hat H_{\rm rec}
=
\hbar\sum_{\boldsymbol\ell}
\delta_{\boldsymbol\ell}
|\boldsymbol\ell\rangle\langle\boldsymbol\ell|.
\]
展开
\[
|\Psi_{\mathbf n}(t)\rangle
=
\sum_{\boldsymbol\ell}
c_{\boldsymbol\ell}^{(\mathbf n)}(t)
|\boldsymbol\ell;\mathbf n\rangle
\]
并代入
\[
\ii\hbar\partial_t|\Psi_{\mathbf n}\rangle
=
(\hat H_{\rm rec}+\hat H_{\rm Q}^{(M)})
|\Psi_{\mathbf n}\rangle.
\]
左乘 $\langle\boldsymbol\ell;\mathbf n|$ 后，来自相邻格点的两类项分别为
\[
+\ii\hbar
\upsilon_\nu
\sqrt{n_\nu-\ell_\nu+1}\,
c_{\boldsymbol\ell-\mathbf e_\nu}^{(\mathbf n)}
\]
和
\[
-\ii\hbar
\upsilon_\nu^*
\sqrt{n_\nu-\ell_\nu}\,
c_{\boldsymbol\ell+\mathbf e_\nu}^{(\mathbf n)}.
\]
与对角反冲项合并并除以 $\hbar$，得到正文\eqnrefs{eq:multimode_qpinem_lattice_dp47}。

最后区分电子标签与联合标签。若存在
$\boldsymbol\ell\neq\boldsymbol\ell\prime$ 却有
\[
\sum_\nu(\ell_\nu-\ell_\nu\prime)\boldsymbol\kappa_\nu=0,
\]
则电子动量态相同，但固定初始光子记录 $\mathbf n$ 时，联合内积仍含
\[
\prod_\nu\langle n_\nu-\ell_\nu\prime|n_\nu-\ell_\nu\rangle
=\prod_\nu\delta_{\ell_\nu,\ell_\nu\prime}.
\]
所以这些联合基态仍正交，不能只因电子动量重合便相干合并不同光子末记录。
电子边带的概率对未读取的光子记录求和；只有电子及光子末记录都相同时才相干合并路径。
在经典数值场极限中不再保存可区分的光子记录，公度多色的各整数路径便按下一节的振幅卷积合并。
\end{proof}


对多模相干态
$|\{\alpha_\nu\}\rangle$，经典场极限应逐模取
\[
|\alpha_\nu|\to\infty,
\qquad
g_{\mathrm Q,\nu}\to0,
\qquad
\beta_\nu
\equiv
\alpha_\nu g_{\\mathrm Q,\nu}
\quad\text{保持有限}.
\]
若在相应工作窗中各电子平移算符可用完整连续谱上的平移表示，它们彼此对易，电子得到
\begin{equation}
\boxed{
 S_{\rm cl}^{(M)}
 =
 \exp\!\left[
 \sum_\nu
 \left(
 \beta_\nu\hat b_\nu^\dagger
 -\beta_\nu^*\hat b_\nu
 \right)
 \right].}
\label{eq:multicolor_classical_scattering_dp47}
\end{equation}
若各动量步长进一步满足
$\kappa_\nu=m_\nu\kappa_0$，不同颜色的整数路径可能到达同一个最终电子格点。把每个颜色的 Jacobi--Anger 展开相乘，并对所有满足
$\ell=\sum_\nu m_\nu r_\nu$
的路径作相干求和，得到
\begin{equation}
\boxed{
 A_\ell
 =\sum_{\{r_\nu\}}
 \delta_{\ell,\sum_\nu m_\nu r_\nu}
 \prod_\nu
 J_{r_\nu}(2|\beta_\nu|)
 \ee^{\ii r_\nu\varphi_\nu},
 \qquad
 \varphi_\nu=\arg\beta_\nu .}
\label{eq:multicolor_pinem_convolution_dp47}
\end{equation}

\begin{proposition}[可交换多色平移的 Bessel 卷积]
设有限个单位平移 $T_\nu$ 及其逆彼此对易，$\beta_\nu$ 固定。则散射指数可逐模因子化，每一因子有绝对收敛的 Bessel 整数平移展开。若所有动量步长 $\kappa_\nu=m_\nu\kappa_0$ 公度，作用到同一电子末态的全部路径相干求和，给\eqnrefs{eq:multicolor_pinem_convolution_dp47}。
\end{proposition}
\begin{proof}
在多模相干态经典极限中，每个量子模式的算符被分解成平均幅与量子涨落；固定
\[
\beta_\nu=\alpha_\nu g_{\\mathrm Q,\nu}
\]
并令相对涨落修正消失后，电子散射算符为
\[
S_{\rm cl}^{(M)}
=
\exp\!\left[
\sum_\nu
(\beta_\nu\hat b_\nu^\dagger-\beta_\nu^*\hat b_\nu)
\right].
\]
完整连续动量空间中的平移算符彼此对易，
\[
[\hat b_\mu,\hat b_\nu]=
[\hat b_\mu,\hat b_\nu^\dagger]=0,
\]
所以指数可逐模分解：
\[
S_{\rm cl}^{(M)}
=
\prod_\nu
\exp(
\beta_\nu\hat b_\nu^\dagger
-\beta_\nu^*\hat b_\nu).
\]

对单个模式，记
\[
\beta_\nu=|\beta_\nu|\ee^{\ii\varphi_\nu},
\qquad
\hat T_\nu\equiv\hat b_\nu^\dagger.
\]
在理想双边平移空间中 $\hat T_\nu$ 是幺正算符，
\[
\hat T_\nu^{-1}=\hat b_\nu.
\]
因此任意整数幂统一写作
\[
\hat T_\nu^{r}
=
\begin{cases}
(\hat b_\nu^\dagger)^r,&r\ge0,\\
(\hat b_\nu)^{-r},&r<0.
\end{cases}
\]
取 $\hat T_\nu$ 的相位本征态
\[
\hat T_\nu|\theta_\nu\rangle
=
\ee^{-\ii\theta_\nu}|\theta_\nu\rangle,
\]
则单模生成元在该表象中成为普通复数：
\begin{align*}
\beta_\nu\hat T_\nu-\beta_\nu^*\hat T_\nu^{-1}
&\longrightarrow
|\beta_\nu|
\left[
\ee^{\ii(\varphi_\nu-\theta_\nu)}
-\ee^{-\ii(\varphi_\nu-\theta_\nu)}
\right]\\
&=
2\ii|\beta_\nu|
\sin(\varphi_\nu-\theta_\nu).
\end{align*}
Jacobi--Anger 恒等式给
\[
\exp[
2\ii|\beta_\nu|\sin(\varphi_\nu-\theta_\nu)
]
=
\sum_{r_\nu\in\mathbb Z}
J_{r_\nu}(2|\beta_\nu|)
\ee^{\ii r_\nu\varphi_\nu}
\ee^{-\ii r_\nu\theta_\nu}.
\]
把
$\ee^{-\ii r_\nu\theta_\nu}$
重新提升为算符 $\hat T_\nu^{r_\nu}$，得到严格的整数平移展开
\[
\exp(
\beta_\nu\hat b_\nu^\dagger
-\beta_\nu^*\hat b_\nu)
=
\sum_{r_\nu\in\mathbb Z}
J_{r_\nu}(2|\beta_\nu|)
\ee^{\ii r_\nu\varphi_\nu}
\hat T_\nu^{r_\nu}.
\]

对所有模式相乘：
\[
S_{\rm cl}^{(M)}
=
\sum_{\{r_\nu\}}
\left[
\prod_\nu
J_{r_\nu}(2|\beta_\nu|)
\ee^{\ii r_\nu\varphi_\nu}
\right]
\prod_\nu\hat T_\nu^{r_\nu}.
\]
若
\[
\kappa_\nu=m_\nu\kappa_0,
\qquad
m_\nu\in\mathbb Z,
\]
则
\[
\prod_\nu\hat T_\nu^{r_\nu}
\]
把电子总共平移
\[
\Delta p
=
\hbar\kappa_0
\sum_\nu m_\nu r_\nu.
\]
因此所有满足
\[
\ell=\sum_\nu m_\nu r_\nu
\]
的整数路径都到达同一个电子末态 $|\ell\rangle$，其振幅必须相干求和。于是
\[
A_\ell
=
\sum_{\{r_\nu\}}
\delta_{\ell,\sum_\nu m_\nu r_\nu}
\prod_\nu
J_{r_\nu}(2|\beta_\nu|)
\ee^{\ii r_\nu\varphi_\nu},
\]
即正文\eqnrefs{eq:multicolor_pinem_convolution_dp47}。

若不同颜色的动量步长不满足这种整数公度关系，则一般不存在唯一的一维总阶数 $\ell$；此时末态应继续用真实电子动量或多指标通道标记，不能把不同可分辨末态强行卷积到一条 Bessel 能量格。
\end{proof}

因此多色 PINEM 的相位敏感性来自“不同交换路径汇聚到同一电子末态”后的振幅干涉；若路径对应不同可分辨电子动量，则应保留多指标末态，而不能先把概率卷积到一条能量格上。


\subsection{单模交换哈密顿量}

现在取一个谱上孤立的光模，假定在实际作用时间内，其余模式及反旋过程对输出态的贡献均低于指定精度；同时假定实际占据电子通道的逐边耦合与失谐可以用共同值代替。这些模型条件的谱隔离与反冲检验将在\chapref{chap:electron-recoil-channels}中分别进行。电子通道因此可近似编号为 $\ell\in\mathbb Z$。定义电子平移算符与光子湮灭算符的作用为
\[
\hat b^\dagger|\ell\rangle_e=|\ell+1\rangle_e,\qquad
\hat b|\ell\rangle_e=|\ell-1\rangle_e,\qquad
\hat a|n\rangle_{\rm ph}=\sqrt n\,|n-1\rangle_{\rm ph}.
\]
前两者改变电子通道编号，后者改变光子数。它们作用在不同的态空间上，因而彼此对易。电子吸收一个光子时，光子数减少一而电子阶数增加一，对应乘积 $\hat b^\dagger\hat a$；逆过程对应 $\hat b\hat a^\dagger$。

沿用真实反冲格的耦合率 $\mathcal G_0(t)$ 和共同失谐 $\Delta_0$，把相互作用绘景中的速率记为
\[
\upsilon(t)\equiv-\ii\mathcal G_0(t)\ee^{\ii\Delta_0(t-t_i)},
\qquad [\upsilon]={\rm s^{-1}}.
\]
因子 $-\ii$ 是相位约定，使传播子中的累积耦合可直接写为
$g_{\mathrm Q}=\int_{t_i}^{t_f}\dd t\,\upsilon(t)$。$g_{\mathrm Q}$ 无量纲。若光场处在振幅为 $\alpha_{\rm coh}(t)$ 的相干态，经典耦合为 $\beta=\int\dd t\,\alpha_{\rm coh}(t)\upsilon(t)$；相干幅度为常数时即 $\beta=\alpha_{\rm coh}g_{\mathrm Q}$。若以后把电子限制为一条二态交换边，常耦合记号 $g_c=\ii\upsilon$ 仍具有频率单位。由上述吸收与发射过程，联合 Hamiltonian 写成
\begin{equation}
\boxed{
\hat H_Q(t)=\ii\hbar[\upsilon(t)\hat b^\dagger\hat a-\upsilon^*(t)\hat b \hat a^\dagger].}
\label{eq:qpinem-single-mode-hamiltonian-dp68}
\end{equation}
其中 $[\upsilon]={\rm s^{-1}}$。共同失谐已经吸收到 $\upsilon$ 的相位中，当前方程便不再另加 $\hbar\Delta_0\Lambda_e$。电子双边格只是局部近似，而光子数仍有 $n=0$ 下界，因此量子边界误差比经典情况多出玻色增强因子。

若 $\Pi_{\min}$ 投影到电子最低格点附近，沿单边近似演化定义的量子边界传播误差为
\begin{equation}
\boxed{
\epsilon_{\rm edge,Q}^{\rm dyn}
=\int_{t_i}^{t_f}\dd t\,|\upsilon(t)|
\sqrt{\langle\Pi_{\min}(\hat n+1)\rangle_t}\ll1.}
\label{eq:quantum_bilateral_ladder_dynamical_error}
\end{equation}
若实际光子数扇区截断在 $n\le N$，则
$\epsilon_{\rm edge,Q}^{\rm dyn}\le\Lambda_I\sqrt{(N+1)P_{\rm edge}}$。因此经典与量子双边延拓不能共用同一个边界判据。

\begin{proposition}[物理格边界的量子传播误差]
设单边电子格嵌入双边格，单边参考 Hamiltonian 在物理子空间外作自伴直和延拓，两个传播子存在。若参考态的边界光子矩可积，则 $\|(U_b-U_s)\Psi_i\|\le\int|\upsilon(t)|\sqrt{\langle\Pi_{\min}(n+1)\rangle_t}\dd t$，即\eqnrefs{eq:quantum_bilateral_ladder_dynamical_error}。只有参考演化始终位于 $n\le N$ 时，才能进一步使用有限 $N$ 上界。
\end{proposition}
\begin{proof}
真实电子能量格存在最低格点 $|\ell_{\min}\rangle$。 令 $s$ 表示物理单边平移算符、$\hat b$ 表示双边延拓后的幺正平移算符，则

\begin{align*}
s|\ell\rangle&=|\ell-1\rangle,
&&\ell>\ell_{\min},\\
s|\ell_{\min}\rangle&=0,\\
\hat b|\ell\rangle&=|\ell-1\rangle,
&&\ell\in\mathbb Z.
\end{align*}
因此在物理单边子空间内
\begin{equation*}
(\hat b-s)\Pi_{\rm phys}
=|\ell_{\min}-1\rangle\langle\ell_{\min}|,
\qquad
(\hat b^\dagger-s^\dagger)\Pi_{\rm phys}=0.
\end{equation*}
双边模型与单边模型的交换哈密顿量分别记为 $H_b(t)$ 与 $H_s(t)$。由
\begin{equation*}
H_Q(t)=\ii\hbar[\upsilon(t)\hat b^\dagger\hat a-\upsilon^*(t)\hat b\hat a^\dagger]
\end{equation*}
可得它们在物理子空间上的差
\begin{align*}
R_{\rm edge}(t)
&\equiv[H_b(t)-H_s(t)]\Pi_{\rm phys}\\
&=-\ii\hbar\upsilon^*(t)
|\ell_{\min}-1\rangle\langle\ell_{\min}|\,\hat a^\dagger .
\end{align*}
整体相位符号不影响范数。对单边参考演化态 $|\Psi_s(t)\rangle$，令
\begin{equation*}
\Pi_{\min}=|\ell_{\min}\rangle\langle\ell_{\min}|,
\qquad
\hat n=\hat a^\dagger\hat a.
\end{equation*}
则
\begin{align*}
\|R_{\rm edge}(t)|\Psi_s(t)\rangle\|^2
&=\hbar^2|\upsilon(t)|^2
\langle\Psi_s|
\Pi_{\min}\hat a\hat a^\dagger\Pi_{\min}
|\Psi_s\rangle\\
&=\hbar^2|\upsilon(t)|^2
\langle\Pi_{\min}(\hat n+1)\rangle_t,
\end{align*}
其中第二步使用 $\hat a\hat a^\dagger=\hat n+1$，并利用电子投影与光子算符作用在不同 Hilbert 空间而彼此对易。两种传播子满足 Duhamel 恒等式
\begin{equation*}
U_b(t_f,t_i)-U_s(t_f,t_i)
=-\frac{\ii}{\hbar}
\int_{t_i}^{t_f}\dd t\,
U_b(t_f,t)R_{\rm edge}(t)U_s(t,t_i).
\end{equation*}
左右传播子均为幺正算符，因此对给定初态取范数并使用三角不等式得到
\begin{align*}
\|(U_b-U_s)|\Psi_i\rangle\|
&\le
\frac1\hbar\int_{t_i}^{t_f}\dd t\,
\|R_{\rm edge}(t)|\Psi_s(t)\rangle\|\\
&=
\int_{t_i}^{t_f}\dd t\,|\upsilon(t)|
\sqrt{\langle\Pi_{\min}(\hat n+1)\rangle_t}.
\end{align*}
这就是\eqnrefs{eq:quantum_bilateral_ladder_dynamical_error}。若参考演化始终限制在 $0\le n\le N$，则算符不等式
\begin{equation*}
\Pi_{\min}(\hat n+1)
\le (N+1)\Pi_{\min}
\end{equation*}
给出
\begin{equation*}
\epsilon_{\rm edge,Q}^{\rm dyn}
\le
\left(\int_{t_i}^{t_f}\dd t\,|\upsilon(t)|\right)
\sqrt{(N+1)P_{\rm edge}},
\end{equation*}
即正文给出的有限光子数上界。
\end{proof}


单模量子交换哈密顿量同时依赖环境单模隔离、电子单动量步长投影、格点平移对称和旋波近似（RWA）。前章\eqnrefs{eq:feqo-fast-branch-propagator-bound}给出了有界工作模型中反旋支路的有限时间充分误差界；谱上孤立或一阶快振荡相消都不自动保证虚过程相位很小。任一误差超出目标精度时，需保留相应被舍通道。

\subsection{Magnus 闭合与守恒量}

含时 Hamiltonian 在不同时刻一般不对易，因此传播子需要保留时间排序。为了求解析表达式，可把传播子写成一个指数 $\exp(\Omega_1+\Omega_2+\cdots)$：第一项 $\Omega_1$ 是相互作用的时间积分，第二项 $\Omega_2$ 修正不同时刻的作用次序，更高项含嵌套对易子。这称为 Magnus 展开。

对\eqnrefs{eq:qpinem-single-mode-hamiltonian-dp68}，电子双边平移满足 $\hat b^\dagger\hat b=\hat b\hat b^\dagger=I$，光子算符满足 $[\hat a,\hat a^\dagger]=I$。两时刻 Hamiltonian 的对易子于是正比于恒等算符，再与 Hamiltonian 对易就为零。因此本模型只需保留前两项。把对应时间积分分别记为累积耦合与整体相位，
\begin{equation}
\boxed{
g_{\mathrm Q}=\int_{t_i}^{t_f}\dd t\,\upsilon(t),
\qquad
\Phi_M
=\frac{1}{2\ii}
\int_{t_i}^{t_f}\dd t_1
\int_{t_i}^{t_1}\dd t_2\,
\bigl[\upsilon^*(t_1)\upsilon(t_2)-\upsilon(t_1)\upsilon^*(t_2)\bigr].}
\label{eq:qpinem-magnus-pulse-parameters-dp68}
\end{equation}
在\eqnrefs{eq:qpinem-single-mode-hamiltonian-dp68} 的模型定义域内，完整时间有序传播子为
\begin{equation}
\boxed{
\hat U_Q(t_f,t_i)
=\ee^{\ii\Phi_M}
\exp[g_{\mathrm Q}\,\hat b^\dagger\hat a-g_{\mathrm Q}^*\hat b\hat a^\dagger].}
\label{eq:qpinem-magnus-exact-dp68}
\end{equation}
$\ee^{\ii\Phi_M}$ 是与电子和光子量子数都无关的整体相位，因此所有单次概率和约化密度算符都不受它影响。若 $\upsilon(t)$ 在整个脉冲中具有固定复相位，则被积的反对称组合恒为零，$\Phi_M=0$；此时正文常用的交换算符可简写为
\begin{equation}
\boxed{
S_Q(g_{\mathrm Q})=\exp[g_{\mathrm Q}\,\hat b^\dagger\hat a-g_{\mathrm Q}^*\hat b\hat a^\dagger].}
\label{eq:qpinem_S}
\end{equation}
电子平移算符与光子升降算符分属不同的 Hilbert 空间；$\hat b$ 是幺正平移算符，而 $\hat a$ 是玻色湮灭算符。两者分别实现电子平移代数与玻色 Fock 代数。

\begin{theorem}[复合交换代数的精确 Magnus 终止]
设双边电子平移与玻色模作用在不同张量因子上，$B=b^\dagger a$。在有限粒子公共核上 $[B,B^\dagger]=I$；若包络可积，对应 Weyl 传播存在，则时间有序指数精确为\eqnrefs{eq:qpinem-magnus-exact-dp68}。二阶项为实整体相位，更高嵌套交换子为零；此结论不适用于破坏双边平移代数的边界截断。
\end{theorem}
\begin{proof}
先把“时间有序指数为什么能够求闭式”化成一个纯算符代数问题。定义
\[
\hat B\equiv\hat b^\dagger\hat a,
\qquad
\hat B^\dagger=\hat b\hat a^\dagger.
\]
电子平移算符与光子算符作用在张量积 Hilbert 空间的不同因子上，故彼此对易；理想双边电子格又满足 $\hat b^\dagger\hat b=\hat b\hat b^\dagger=I_e$。于是
\begin{align*}
\hat B\hat B^\dagger
&=\hat b^\dagger\hat b\,\hat a\hat a^\dagger
=I_e(\hat N_{\rm ph}+1),\\
\hat B^\dagger\hat B
&=\hat b\hat b^\dagger\,\hat a^\dagger\hat a
=I_e\hat N_{\rm ph},
\end{align*}
从而
\begin{equation*}
[\hat B,\hat B^\dagger]=I_e\otimes I_{\rm ph}\equiv I.
\end{equation*}
因此复合交换算符 $\hat B$、$\hat B^\dagger$ 在这个有效 Hilbert 空间中形成与单个玻色模相同的 Heisenberg--Weyl 对易代数；这正是闭式存在的根源。

传播子满足 $\dot U=\hat A(t)U$，其中
\begin{equation*}
\hat A(t)=-\frac{\ii}{\hbar}\hat H_Q(t)
=\upsilon(t)\hat B-\upsilon^*(t)\hat B^\dagger.
\end{equation*}
Magnus 展开定义
\begin{align*}
\Omega_1&=\int_{t_i}^{t_f}\dd t_1\,\hat A(t_1),\\
\Omega_2&=\frac12\int_{t_i}^{t_f}\dd t_1
\int_{t_i}^{t_1}\dd t_2\,[\hat A(t_1),\hat A(t_2)],
\end{align*}
而 $\Omega_3$ 及更高项只含形如 $[\hat A(t_1),[\hat A(t_2),\hat A(t_3)]]$ 的嵌套对易子。先显式计算二时刻对易子：
\begin{align*}
[\hat A_1,\hat A_2]
={}&-\upsilon_1\upsilon_2^*[\hat B,\hat B^\dagger]
-\upsilon_1^*\upsilon_2[\hat B^\dagger,\hat B]\\
={}&\upsilon_1^*\upsilon_2-\upsilon_1\upsilon_2^*.
\end{align*}
右端乘的是单位算符，所以对任意 $t_3$ 都有
\[
[\hat A(t_3),[\hat A(t_1),\hat A(t_2)]]=0.
\]
所有含三重及更高嵌套对易子的 Magnus 项因此同时消失，而不是逐阶近似舍去。

第一项直接积分为
\begin{align*}
\Omega_1
&=\left(\int_{t_i}^{t_f}\dd t\,\upsilon(t)\right)\hat B
-\left(\int_{t_i}^{t_f}\dd t\,\upsilon^*(t)\right)\hat B^\dagger\\
&=g_{\mathrm Q}\hat B-g_{\mathrm Q}^*\hat B^\dagger.
\end{align*}
第二项则为
\begin{align*}
\Omega_2
&=\frac12\int_{t_i}^{t_f}\dd t_1
\int_{t_i}^{t_1}\dd t_2\,
[\upsilon^*(t_1)\upsilon(t_2)-\upsilon(t_1)\upsilon^*(t_2)]I\\
&=\ii\Phi_M I.
\end{align*}
被积括号等于 $2\ii\,\operatorname{Im}[\upsilon^*(t_1)\upsilon(t_2)]$，所以 $\Phi_M$ 为实数，$\ee^{\ii\Phi_M}$ 确为纯整体相位。由于 $[\Omega_1,\Omega_2]=0$，
\begin{align*}
U_Q(t_f,t_i)
&=\exp(\Omega_1+\Omega_2)\\
&=\ee^{\ii\Phi_M}
\exp[g_{\mathrm Q}\hat b^\dagger\hat a-g_{\mathrm Q}^*\hat b\hat a^\dagger],
\end{align*}
即正文\eqnrefs{eq:qpinem-magnus-exact-dp68}。
\end{proof}


在同一单模双边交换模型内，生成元具有守恒量
\begin{equation}
\boxed{
\mathcal N_{\rm ex}=\Lambda_e+N_{\rm ph},
\qquad
\Lambda_e=\sum_{\ell\in\mathbb Z}\ell|\ell\rangle\langle\ell|,
\qquad
N_{\rm ph}=\hat a^\dagger\hat a.}
\label{eq:qpinem-exchange-number-dp68}
\end{equation}
因为每次电子上移一个能量步长都伴随一个光子被湮灭，反过程则伴随光子产生。于是
\begin{equation}
\boxed{[\mathcal N_{\rm ex},S_Q]=0.}
\label{eq:qpinem_conservation}
\end{equation}
若初态为 $|0,n\rangle$，所有末态必须满足 $\ell+m=n$。这条选择规则来自交换数守恒，与对 $g_{\mathrm Q}$ 的微扰展开无关；在单模交换 Hamiltonian 的适用域内，对任意有限 $g_{\mathrm Q}$ 都成立。

\begin{proposition}[联合交换数的守恒与谱扇区]
在双边电子格与玻色 Fock 空间上，$\mathcal N_{\rm ex}=\Lambda_e+N_{\rm ph}$ 与两个交换顶角对易。对存在的 Weyl 传播，守恒数的每个谱子空间均保持不变，$[\mathcal N_{\rm ex},U_Q]=0$；因此 $|0,n\rangle$ 只连接 $|\ell,m\rangle$ 且 $\ell+m=n$。
\end{proposition}
\begin{proof}
先从电子格基底本身推出平移算符的对易关系，而不把它当作已知公式。由
\[
\Lambda_e|\ell\rangle=\ell|\ell\rangle,
\qquad
\hat b^\dagger|\ell\rangle=|\ell+1\rangle,
\qquad
\hat b|\ell\rangle=|\ell-1\rangle,
\]
有
\begin{align*}
[\Lambda_e,\hat b^\dagger]|\ell\rangle
&=(\ell+1)|\ell+1\rangle-\ell|\ell+1\rangle
=\hat b^\dagger|\ell\rangle,\\
[\Lambda_e,\hat b]|\ell\rangle
&=(\ell-1)|\ell-1\rangle-\ell|\ell-1\rangle
=-\hat b|\ell\rangle.
\end{align*}
因为这对所有整数 $\ell$ 成立，故
\[
[\Lambda_e,\hat b^\dagger]=\hat b^\dagger,
\qquad [\Lambda_e,\hat b]=-\hat b.
\]
同理，由 $N_{\rm ph}|n\rangle=n|n\rangle$ 与玻色升降关系得到
\[
[N_{\rm ph},\hat a]=-\hat a,
\qquad
[N_{\rm ph},\hat a^\dagger]=\hat a^\dagger.
\]
于是联合交换数
\[
\mathcal N_{\rm ex}=\Lambda_e+N_{\rm ph}
\]
对两个基本交换顶角分别给出
\begin{align*}
[\mathcal N_{\rm ex},\hat b^\dagger\hat a]
&=[\Lambda_e,\hat b^\dagger]\hat a
+\hat b^\dagger[N_{\rm ph},\hat a]
=0,\\
[\mathcal N_{\rm ex},\hat b\hat a^\dagger]
&=[\Lambda_e,\hat b]\hat a^\dagger
+\hat b[N_{\rm ph},\hat a^\dagger]
=0.
\end{align*}
因此对任意时变包络 $\upsilon(t)$ 都有
\[
[\mathcal N_{\rm ex},\hat H_Q(t)]=0.
\]
为了把“瞬时对易”严格传到有限时间传播子，定义 Heisenberg 型量
\[
\mathcal N_U(t)=U_Q^\dagger(t,t_i)\mathcal N_{\rm ex}U_Q(t,t_i).
\]
利用
$\ii\hbar\,\partial_tU_Q=\hat H_QU_Q$
和
$\partial_tU_Q^\dagger=(\ii/\hbar)U_Q^\dagger\hat H_Q$
可得
\begin{align*}
\frac{\dd\mathcal N_U}{\dd t}
&=\frac{\ii}{\hbar}U_Q^\dagger\hat H_Q\mathcal N_{\rm ex}U_Q
-\frac{\ii}{\hbar}U_Q^\dagger\mathcal N_{\rm ex}\hat H_QU_Q\\
&=\frac{\ii}{\hbar}U_Q^\dagger[\hat H_Q,\mathcal N_{\rm ex}]U_Q=0.
\end{align*}
所以 $\mathcal N_U(t)=\mathcal N_{\rm ex}$，等价于
\[
[\mathcal N_{\rm ex},U_Q(t_f,t_i)]=0.
\]
固定相位时 $U_Q$ 除整体相位外就是 $S_Q$，从而得到正文\eqnrefs{eq:qpinem_conservation}。

守恒扇区也可直接从谱投影看出。若初态 $|0,n\rangle$ 满足
$\mathcal N_{\rm ex}|0,n\rangle=n|0,n\rangle$，则
\begin{align*}
\mathcal N_{\rm ex}U_Q|0,n\rangle
=U_Q\mathcal N_{\rm ex}|0,n\rangle
=nU_Q|0,n\rangle.
\end{align*}
所以输出态仍完全位于本征值为 $n$ 的子空间；基底态 $|\ell,m\rangle$ 的必要条件为 $\ell+m=n$。
\end{proof}


若 $\upsilon(t)=\ee^{\ii\varphi_0}u(t)$ 且 $u(t)$ 为实包络，则 $\upsilon^*(t_1)\upsilon(t_2)$ 为实数，从而 $\Phi_M=0$；此时传播子退化为正文常用的交换算符 $S_Q(g_{\mathrm Q})$。

\subsection{Fock 响应}

守恒量已经把联合态限制在一条 $\ell+m=n$ 的格线上。对初始电子中心格点和光子 Fock 态 $|n\rangle$，该有效模型内的完整输出写成
\begin{equation}
\boxed{
|0,n\rangle\xrightarrow{S_Q}
\sum_{\ell=-\infty}^{n}
\mathcal A_{\ell|n}(g_{\mathrm Q})|\ell,n-\ell\rangle.}
\label{eq:qpinem_joint_fock}
\end{equation}
由于电子相位基底把 $\hat b$ 对角化，$S_Q$ 在每个电子相位扇区上等价于一个光子位移算符。结果可写成单模位移算符矩阵元
\begin{equation}
\boxed{
\mathcal A_{\ell|n}(g_{\mathrm Q})
=\langle n-\ell|D(-g_{\mathrm Q}^*)|n\rangle.}
\label{eq:A_l_n_displacement}
\end{equation}
把位移算符矩阵元展开后，令 $r=|g_{\mathrm Q}|^2$，可把正文\eqnrefs{eq:A_l_n_displacement} 写成显式的分段闭式
\begin{equation}
\boxed{
\mathcal A_{\ell|n}=\ee^{-r/2}
\begin{cases}
\displaystyle
\sqrt{\frac{(n-\ell)!}{n!}}\,g_{\mathrm Q}^{\ell}
L_{n-\ell}^{(\ell)}(r), & 0\le\ell\le n,\\[8pt]
\displaystyle
\sqrt{\frac{n!}{(n-\ell)!}}\,(-g_{\mathrm Q}^*)^{-\ell}
L_n^{(-\ell)}(r), & \ell<0,
\end{cases}}
\label{eq:qpinem-fock-laguerre-dp68}
\end{equation}
其中 $L_j^{(k)}$ 是缔合 Laguerre 多项式。第一支有上界 $\ell\le n$，因为电子净增益 $\ell$ 个量子必须从初始 Fock 态中湮灭同样数目的光子；电子净损失 $\ell<0$ 则可以任意地产生光子。因而在该单模双边交换模型内，完整 Fock 响应由 Laguerre 多项式控制；Bessel 分布只在相干态经典场极限中恢复。

\begin{lemma}[位移算符的 Fock--Laguerre 矩阵元]
设 $D(\alpha)=\exp(\alpha a^\dagger-\alpha^*a)$。对 $m\ge n$，$\langle m|D(\alpha)|n\rangle=e^{-|\alpha|^2/2}\sqrt{n!/m!}\,\alpha^{m-n}L_n^{(m-n)}(|\alpha|^2)$；对 $m\le n$ 由共轭得到另一支。电子 Fourier 投影后，两支给\eqnrefs{eq:qpinem-fock-laguerre-dp68}，不存在的负光子数末态振幅为零。
\end{lemma}
\begin{proof}
先解释为什么电子平移问题会变成光子的位移算符问题。定义电子相位态
\begin{equation*}
 |\theta\rangle=\frac{1}{\sqrt{2\pi}}
 \sum_{\ell\in\mathbb Z}\ee^{\ii\ell\theta}|\ell\rangle,
 \qquad
 \hat b|\theta\rangle=\ee^{\ii\theta}|\theta\rangle .
\end{equation*}
于是固定 $\theta$ 后，交换算符中的电子算符退化为复数：
\begin{align*}
 S_Q(g_{\mathrm Q})|\theta\rangle|n\rangle
 &=|\theta\rangle
 \exp\!\left[
 g_{\mathrm Q} \ee^{-\ii\theta}\hat a
 -g_{\mathrm Q}^*\ee^{+\ii\theta}\hat a^\dagger
 \right]|n\rangle\\
 &=|\theta\rangle
 D\!\left(-g_{\mathrm Q}^*\ee^{+\ii\theta}\right)|n\rangle ,
\end{align*}
其中 $D(\alpha)=\exp(\alpha\hat a^\dagger-\alpha^*\hat a)$。将末态电子投影到 $|\ell\rangle$，相位积分只提取位移矩阵元中 Fourier 次数为 $\ell$ 的分量，得到正文\eqnrefs{eq:A_l_n_displacement}。真正需要计算的是一般的 $\langle m|D(\alpha)|n\rangle$。

因为 $[\alpha\hat a^\dagger,-\alpha^*\hat a]=|\alpha|^2$ 是 $c$ 数，Baker--Campbell--Hausdorff 公式在这里精确给出
\begin{equation*}
 D(\alpha)
 =\ee^{-|\alpha|^2/2}
 \ee^{\alpha\hat a^\dagger}\ee^{-\alpha^*\hat a}.
\end{equation*}
先让右边的湮灭指数作用在 $|n\rangle$：
\begin{equation*}
 \ee^{-\alpha^*\hat a}|n\rangle
 =\sum_{k=0}^{n}
 \frac{(-\alpha^*)^k}{k!}
 \sqrt{\frac{n!}{(n-k)!}}\,|n-k\rangle .
\end{equation*}
若 $m\ge n$，随后产生指数必须再产生
$j=m-n+k$ 个光子才可能投影到 $\langle m|$，因此
\begin{align*}
 \langle m|D(\alpha)|n\rangle
 &=\ee^{-|\alpha|^2/2}\sqrt{m!n!}
 \sum_{k=0}^{n}
 \frac{(-\alpha^*)^k\alpha^{m-n+k}}
 {k!(m-n+k)!(n-k)!}\\
 &=\ee^{-|\alpha|^2/2}
 \sqrt{\frac{n!}{m!}}\,
 \alpha^{m-n}L_n^{(m-n)}(|\alpha|^2).
\end{align*}
最后一步只用了缔合 Laguerre 多项式的有限和定义
\begin{equation*}
 L_n^{(r)}(x)
 =\sum_{k=0}^{n}
 \frac{(n+r)!}{(n-k)!(r+k)!k!}\,
 (-x)^k.
\end{equation*}
对 $m\le n$ 可重复同一计算，或利用
$D^\dagger(\alpha)=D(-\alpha)$ 取厄米共轭，得到
\begin{equation*}
 \langle m|D(\alpha)|n\rangle
 =\ee^{-|\alpha|^2/2}
 \sqrt{\frac{m!}{n!}}\,
 (-\alpha^*)^{n-m}L_m^{(n-m)}(|\alpha|^2).
\end{equation*}
现在代入 $m=n-\ell$、$\alpha=-g_{\mathrm Q}^*$。当 $0\le\ell\le n$ 时属于 $m\le n$ 支，且 $-\alpha^*=g_{\mathrm Q}$；当 $\ell<0$ 时属于 $m>n$ 支，且 $\alpha=-g_{\mathrm Q}^*$。两支分别给出正文\eqnrefs{eq:qpinem-fock-laguerre-dp68}。
\end{proof}


当孤立玻色模式取为近场光子模式，电子动量步长由该模式的主导纵向 Fourier 分量选出，并且电子双边格误差满足\eqnrefs{eq:quantum_bilateral_ladder_dynamical_error} 时，\eqnrange{eq:qpinem-single-mode-hamiltonian-dp68}{eq:qpinem-fock-laguerre-dp68} 就构成受控的自由电子--单光子模式 Fock 交换模型。它的实验输入是模式频率、归一化场分布、有限作用时间和初始光子统计；模式重叠、逐边反冲可分辨或电子谱边界被占据时，应返回多模连续谱或有限反冲联合格。

弱耦合时，令 $r_{\rm ph}\in\mathbb N_+$ 表示电子净增益阶数；第 $r_{\rm ph}$ 阶增益需要从光场中湮灭 $r_{\rm ph}$ 个光子，故最低阶概率正比于同阶光子数下降阶乘：
\begin{equation}
\boxed{
P_{+r_{\rm ph}}
=\frac{|g_{\mathrm Q}|^{2r_{\rm ph}}}{(r_{\rm ph}!)^2}
\left\langle (\hat a^\dagger)^{r_{\rm ph}} \hat a^{r_{\rm ph}}\right\rangle
+O(|g_{\mathrm Q}|^{2r_{\rm ph}+2}).}
\label{eq:qpinem-weak-factorial-moment-dp19}
\end{equation}

同一展开在 $r_{\rm ph}=1$ 时给
$P_{+1}=|g_{\mathrm Q}|^2\langle \hat a^\dagger\hat a\rangle+O(|g_{\mathrm Q}|^4)$，所以
\begin{equation*}
\frac{(r_{\rm ph}!)^2P_{+r_{\rm ph}}}{P_{+1}^{r_{\rm ph}}}
=\frac{\langle(\hat a^\dagger)^{r_{\rm ph}}\hat a^{r_{\rm ph}}\rangle}
{\langle \hat a^\dagger\hat a\rangle^{r_{\rm ph}}}+O(|g_{\mathrm Q}|^2)
=g_{\rm ph}^{(r_{\rm ph})}(0)+O(|g_{\mathrm Q}|^2).
\end{equation*}
这建立了“电子边带作为光子统计探测器”的概率关系：光子数分布先经条件响应核映射为理想电子边带概率；有限探测窗口、效率与读出噪声随后由测量章定义的 POVM 或量子仪器作用在这些理想边带上。

\begin{proposition}[弱交换的光子阶乘矩]
设电子初态为单一格点，光场初态在有限 Fock 数窗支撑，正整数 $r_{\rm ph}$ 固定。当 $g_{\mathrm Q}\to0$ 时，电子吸收 $r_{\rm ph}$ 个量子的概率为 $|g_{\mathrm Q}|^{2r_{\rm ph}}\langle(a^\dagger)^{r_{\rm ph}}a^{r_{\rm ph}}\rangle/(r_{\rm ph}!)^2+O(|g_{\mathrm Q}|^{2r_{\rm ph}+2})$。无限光子尾的态需额外矩条件和一致余项控制，不能把有限窗的算符范数展开称为任意态的一致展开。
\end{proposition}
\begin{proof}
令
\[
\hat B=\hat b^\dagger\hat a,
\qquad [\hat B,\hat B^\dagger]=I,
\qquad r\equiv|g_{\mathrm Q}|^2.
\]
由 BCH 公式，交换算符可精确正规排序为
\begin{equation*}
S_Q(g_{\mathrm Q})
=\ee^{-r/2}
\exp(-g_{\mathrm Q}^*\hat B^\dagger)
\exp(g_{\mathrm Q}\hat B).
\end{equation*}
这比直接把整个指数按总阶数展开更适合判断“电子净移动了多少格”。对电子初态 $|0\rangle$ 与末态 $|r_{\rm ph}\rangle$，定义只作用在光场上的条件振幅算符
\[
\hat M_{+r_{\rm ph}}
\equiv{}_e\!\langle r_{\rm ph}|S_Q|0\rangle_e.
\]
把两个指数分别展开：
\begin{align*}
\hat M_{+r_{\rm ph}}
=\ee^{-r/2}
\sum_{j,k\ge0}
\frac{(-g_{\mathrm Q}^*)^jg_{\mathrm Q}^k}{j!k!}
{}_e\!\langle r_{\rm ph}|
\hat b^j(\hat b^\dagger)^k|0\rangle_e
(\hat a^\dagger)^j\hat a^k.
\end{align*}
双边格上 $(\hat b^\dagger)^k$ 使电子上移 $k$ 格，而 $\hat b^j$ 再使它下移 $j$ 格，因此电子矩阵元严格给出整数选择条件
\[
k-j=r_{\rm ph}.
\]
令 $k=r_{\rm ph}+j$，得到一个对任意 $g_{\mathrm Q}$ 都成立的条件算符级数
\begin{equation*}
\hat M_{+r_{\rm ph}}
=\ee^{-r/2}
\sum_{j=0}^{\infty}
\frac{(-g_{\mathrm Q}^*)^jg_{\mathrm Q}^{r_{\rm ph}+j}}
{j!(r_{\rm ph}+j)!}
(\hat a^\dagger)^j\hat a^{r_{\rm ph}+j}.
\end{equation*}
最低项 $j=0$ 恰好要求连续湮灭 $r_{\rm ph}$ 个光子；下一项必须同时多一次电子下移与一次电子上移，所以总耦合阶数增加二，而不是一。再把 $\ee^{-r/2}=1+O(|g_{\mathrm Q}|^2)$ 展开，得到
\begin{equation*}
\hat M_{+r_{\rm ph}}
=\frac{g_{\mathrm Q}^{r_{\rm ph}}}{r_{\rm ph}!}\hat a^{r_{\rm ph}}
+O(|g_{\mathrm Q}|^{r_{\rm ph}+2}).
\end{equation*}
因此对任意光场初态 $\rho_{\rm ph}$，
\begin{align*}
P_{+r_{\rm ph}}
&=\Tr\!\left[
\hat M_{+r_{\rm ph}}\rho_{\rm ph}
\hat M_{+r_{\rm ph}}^\dagger\right]\\
&=\frac{|g_{\mathrm Q}|^{2r_{\rm ph}}}{(r_{\rm ph}!)^2}
\Tr\!\left[
\hat a^{r_{\rm ph}}\rho_{\rm ph}
(\hat a^\dagger)^{r_{\rm ph}}\right]
+O(|g_{\mathrm Q}|^{2r_{\rm ph}+2})\\
&=\frac{|g_{\mathrm Q}|^{2r_{\rm ph}}}{(r_{\rm ph}!)^2}
\left\langle(\hat a^\dagger)^{r_{\rm ph}}\hat a^{r_{\rm ph}}\right\rangle
+O(|g_{\mathrm Q}|^{2r_{\rm ph}+2}),
\end{align*}
即正文\eqnrefs{eq:qpinem-weak-factorial-moment-dp19}。
\end{proof}


% -----------------------------------------------------------------------------
% -----------------------------------------------------------------------------

\begin{exbox}[QPINEM 真空损失]
取电子中心格 $|0\rangle$、光子数态 $|n\rangle$ 入射，求弱交换的增益与损失概率，并比较真空与有光输入。量子交换传播子\eqnrefs{eq:qpinem_S}对 Fock 输入的弱交换条件为 $|g_{\mathrm Q}|^2(n+1)\ll1$。

\textbf{解。} 在实际占据范围内满足该条件后，传播子展开为
\begin{equation*}
 S_Q
 =1+g_{\mathrm Q}\hat b^\dagger\hat a
 -g_{\mathrm Q}^*\hat b\hat a^\dagger+O(|g_{\mathrm Q}|^2).
\end{equation*}
对初态 $|0,n\rangle$，第一项保持电子能量不变；第二项湮灭一个光子并把电子推到 $\ell=+1$；第三项产生一个光子并把电子推到 $\ell=-1$。因此
\begin{align}
 P_{+1|n}&=|g_{\mathrm Q}|^2n+O(|g_{\mathrm Q}|^4),\notag\\
 P_{-1|n}&=|g_{\mathrm Q}|^2(n+1)+O(|g_{\mathrm Q}|^4).
 \label{eq:qpinem-gain-loss-n-nplus1-dp65}
\end{align}
这里的 $+1$ 是电子\emph{增益}一个光子能量，$-1$ 是电子\emph{损失}一个光子能量。两者的差
\begin{equation}
 \boxed{P_{-1|n}-P_{+1|n}=|g_{\mathrm Q}|^2+O(|g_{\mathrm Q}|^4)}
 \label{eq:qpinem-spontaneous-difference-dp65}
\end{equation}
与 $n$ 无关，正是玻色对易关系 $[\hat a,\hat a^\dagger]=1$ 留下的真空项。

例如取 $|g_{\mathrm Q}|=0.10$：真空 $n=0$ 时
\begin{equation*}
 P_{+1|0}=0,
 \qquad
 P_{-1|0}\simeq0.010.
\end{equation*}
所以“没有入射光子”并不意味着电子不能损失能量；电子仍可向真空模式自发发射一个量子。对 $n=1$，增益和损失概率分别约为 $0.010$ 与 $0.020$；对 $n=10$，则约为 $0.10$ 与 $0.11$。只有同时维持 $|g_{\mathrm Q}|^2(n+1)\ll1$ 时，增大 $n$ 才使这两条弱交换率的相对差异按 $1/n$ 减小。这只检验增益--损失对称，不能据此把 Fock 输入的电子约化态视为经典相干调制后的纯态。


\medskip

\eqnrefs{eq:qpinem_S} 给出单模 QPINEM 的完整联合散射算符。在 $|g_{\mathrm Q}|\ll1$ 的弱交换极限，散射算符对初态 $|0,n\rangle$ 的一阶振幅只含一个光子升降算符，因此 Fock 梯矩阵元直接给出 $\sqrt n$ 与 $\sqrt{n+1}$ 的玻色增强。概率保留到 $O(|g_{\mathrm Q}|^2)$ 即得到\eqnrefs{eq:qpinem-gain-loss-n-nplus1-dp65}；电子双边平移仍要求实际占据通道远离谱下界。


由指数展开
\begin{equation*}
 S_Q=1+X+O(X^2),
 \qquad
 X=g_{\mathrm Q}\hat b^\dagger\hat a-g_{\mathrm Q}^*\hat b\hat a^\dagger.
\end{equation*}
作用在 $|0,n\rangle$ 上：
\begin{align*}
 \hat b^\dagger\hat a|0,n\rangle
 &=\sqrt n\,|+1,n-1\rangle,\\
 \hat b\hat a^\dagger|0,n\rangle
 &=\sqrt{n+1}\,|-1,n+1\rangle.
\end{align*}
因此
\begin{equation*}
 S_Q|0,n\rangle
 =|0,n\rangle
 +g_{\mathrm Q}\sqrt n\,|+1,n-1\rangle
 -g_{\mathrm Q}^*\sqrt{n+1}\,|-1,n+1\rangle
 +O(|g_{\mathrm Q}|^2).
\end{equation*}
两个末态的电子边带和光子数都互相正交，故到这一阶不存在交叉项。振幅取模平方即得
\begin{equation*}
 P_{+1|n}=|g_{\mathrm Q}|^2n,
 \qquad
 P_{-1|n}=|g_{\mathrm Q}|^2(n+1).
\end{equation*}
相减后只剩 $|g_{\mathrm Q}|^2$。从算符层面看，这个 $1$ 来自
\begin{equation*}
 \langle n|\hat a\hat a^\dagger|n\rangle
 =\langle n|(\hat a^\dagger\hat a+1)|n\rangle=n+1,
\end{equation*}
所以它不能由经典复振幅替代；经典场极限只是在 $n\gg1$ 时把相对真空修正压到 $1/n$。
\end{exbox}



\begin{exbox}[相同平均光子数的不同电子谱]
现在固定三种光场都具有同一平均光子数
\begin{equation*}
 \bar n=\langle\hat n\rangle=2,
\end{equation*}
但分别取 Fock 态 $|2\rangle$、平均数为 2 的相干态和平均数为 2 的单模热态。取 $|g_{\mathrm Q}|=0.10$，求领先的一阶和二阶增益概率，判断电子计数能否区分这些输入。

\textbf{解。} 它们的一阶平均完全相同，因此在弱交换下都有
\begin{equation*}
 P_{+1}\simeq|g_{\mathrm Q}|^2\bar n=0.020.
\end{equation*}
但二阶增益边带由\eqnrefs{eq:qpinem-weak-factorial-moment-dp19} 控制：
\begin{equation*}
 P_{+2}
 \simeq\frac{|g_{\mathrm Q}|^4}{4}
 \langle\hat n(\hat n-1)\rangle.
\end{equation*}
三种光场的结果是
\begin{center}
\begin{tabular}{c|c|c|c}
 光场 & $\langle\hat n(\hat n-1)\rangle$ & $P_{+2}$ & $g^{(2)}(0)$\\\hline
 $|2\rangle$ & 2 & $5.0\times10^{-5}$ & $1/2$\\
 相干态，$|\alpha|^2=2$ & 4 & $1.0\times10^{-4}$ & $1$\\
 单模热态，$\bar n=2$ & 8 & $2.0\times10^{-4}$ & $2$
\end{tabular}
\end{center}
也就是说，三种场的平均光子数完全相同，电子的一阶增益概率也相同，但二阶边带相差四倍。用只包含电子计数的组合
\begin{equation}
 \boxed{
 \frac{4P_{+2}}{P_{+1}^2}
 =g^{(2)}(0)+O(|g_{\mathrm Q}|^2)}
 \label{eq:qpinem-g2-electron-ratio-dp65}
\end{equation}
即可在弱耦合极限中区分亚 Poisson、Poisson 和聚束光子统计。高阶电子边带把光场的下降阶乘矩逐阶映射到电子能谱，因此相同平均光子数仍可给出不同的二阶增益。
\end{exbox}


\subsection{从位移矩阵元到条件光子数}
前面的 Laguerre 闭式给出全阶单模散射。为把每个阶乘和相位直接对应到跃迁路径，
令 $\zeta=-g_{\mathrm Q}^*$、$\mu=|\zeta|^2$，在固定总交换数的链上使用普通振子位移 $D(\zeta)$。
\begin{proposition}[有限和形式的位移振幅]
对 $m,n\ge0$，
\begin{equation}
 D_{mn}(\zeta)=\ee^{-\mu/2}\sqrt{m!n!}
 \sum_{k=0}^{\min(m,n)}
 \frac{\zeta^{m-k}(-\zeta^*)^{n-k}}{k!(m-k)!(n-k)!}.
 \label{eq:feqo-displacement-matrix}
\end{equation}
入射 $|0,n\rangle$ 的输出为 $\sum_mD_{mn}|n-m,m\rangle$。
\end{proposition}
\begin{proof}
因 $[\zeta a^\dagger,-\zeta^*a]=\mu I$，正规排序给
$D(\zeta)=\ee^{-\mu/2}\ee^{\zeta a^\dagger}\ee^{-\zeta^*a}$。
右侧指数从 $|n\rangle$ 降到 $|k\rangle$，系数为
$(-\zeta^*)^{n-k}\sqrt{n!/k!}/(n-k)!$；左侧再升到 $|m\rangle$，
系数为 $\zeta^{m-k}\sqrt{m!/k!}/(m-k)!$。
对允许的 $0\le k\le\min(m,n)$ 求和即得。
位移算符酉性还给 $\sum_m|D_{mn}|^2=1$，故概率由全部末态承接，不能任意归一化有限的三个边带。
\end{proof}

\begin{exbox}[单光子输入的中心峰相消]
取电子单格与光子数态 $n=1$ 入射，求中心峰和相邻两条边带的全阶振幅，判断中心峰是否总保留非零权重。

\textbf{解。} 在\eqnrefs{eq:feqo-displacement-matrix}中代入 $n=1$，有限和至多含两条路径：
\begin{equation}
 D_{0,1}=-\zeta^*\ee^{-\mu/2},\qquad
 D_{1,1}=(1-\mu)\ee^{-\mu/2},\qquad
 D_{2,1}=\frac{\zeta(2-\mu)}{\sqrt2}\ee^{-\mu/2}.
 \label{eq:feqo-single-photon-exact-amplitudes}
\end{equation}
所以 $P_{+1}=\mu\ee^{-\mu}$、$P_0=(1-\mu)^2\ee^{-\mu}$、
$P_{-1}=\mu(2-\mu)^2\ee^{-\mu}/2$。
在 $\mu=1$ 时无交换峰为零：这是振幅中不交换路径与返回路径相消，
不是光子不存在；概率仍由其他所有边带承担。
\end{exbox}

\begin{definition}[电子结果对应的条件光态]
若 $|\Psi\rangle=\sum_{\ell,m}C_{\ell m}|\ell,m\rangle$，理想电子结果 $\ell$ 的概率与条件态为
\begin{equation}
 p_\ell=\sum_m|C_{\ell m}|^2,\qquad
 |\gamma_\ell\rangle=p_\ell^{-1/2}\sum_mC_{\ell m}|m\rangle,\qquad p_\ell>0.
 \label{eq:feqo-conditional-photon-state}
\end{equation}
条件光子数分布为
\begin{equation}
 P(m|\ell)=|C_{\ell m}|^2/p_\ell.
 \label{eq:feqo-conditional-photon-number}
\end{equation}
零概率结果不定义归一化条件态；有限能窗和混合输入应使用测量章的条件量子仪器。
\end{definition}

\begin{exbox}[相干输入的精确条件分布及经典极限的限制]
取电子单格入射，光场为相干态 $|\alpha\rangle$。给定电子末阶数 $\ell$，求条件光子数分布，并判断它何时接近输入的 Poisson 分布。

\textbf{解。} 输入相干态的数态系数为 $c_n=\ee^{-|\alpha|^2/2}\alpha^n/\sqrt{n!}$。
守恒量要求同一末态 $(\ell,m)$ 只能由初态光子数 $m+\ell$ 到达，故
\begin{equation}
 C_{\ell m}=c_{m+\ell}D_{m,m+\ell}(\zeta),\qquad
 (\rho_e)_{\ell\ell'}=\sum_mC_{\ell m}C_{\ell'm}^*.
 \label{eq:feqo-coherent-reduced-exact}
\end{equation}
负指标 $c_n$ 取零。于是
\begin{equation}
 P(m|\ell)=\frac{|c_{m+\ell}D_{m,m+\ell}(\zeta)|^2}
 {\sum_{m'}|c_{m'+\ell}D_{m',m'+\ell}(\zeta)|^2}.
 \label{eq:feqo-pinem-conditional}
\end{equation}
一般有限耦合下，此分布没有理由是均值 $|\alpha|^2-\ell$ 的 Poisson 分布：
电子后选择还乘上依赖 $m$ 的位移矩阵元。
经典缩放 $g_{\mathrm Q}\to0$、$|\alpha|\to\infty$ 且 $g_{\mathrm Q}\alpha=\beta$ 固定时，
对具有非零极限概率的边带，条件光态趋近输入的宏观相干场，条件数分布与输入 Poisson 数分布趋近。
但若某条边带接近零概率，相同的无条件误差不能直接控制其归一化条件态，
误差会被 $1/p_\ell$ 放大。真空输入则完全不同：结果 $\ell=-m$ 精确预告光子数态 $|m\rangle$，
其 $g^{(2)}(0)=1-1/m$ 仅在 $m\ge1$ 时有定义；$m=0$ 的分母为零，不能称其 $g^{(2)}=0$。
\end{exbox}

\begin{proposition}[弱量子交换读取的相干与数统计]
对初始电子 $|0\rangle$，展开 $S=I+g_{\mathrm Q}b^\dagger a-g_{\mathrm Q}^*ba^\dagger+O(|g_{\mathrm Q}|^2)$。
其中电子增益振幅为 $g_{\mathrm Q}a$，电子损失振幅为 $-g_{\mathrm Q}^*a^\dagger$，因此
\begin{equation}
 P_{+1}=|g_{\mathrm Q}|^2\bar n+O(|g_{\mathrm Q}|^4),\quad
 P_{-1}=|g_{\mathrm Q}|^2(\bar n+1)+O(|g_{\mathrm Q}|^4),\quad
 P_0=1-|g_{\mathrm Q}|^2(2\bar n+1)+O(|g_{\mathrm Q}|^4).
 \label{eq:feqo-weak-quantum-populations}
\end{equation}
此外
\begin{equation}
 (\rho_e)_{1,0}=g_{\mathrm Q}\langle a\rangle+O(|g_{\mathrm Q}|^2),\qquad
 (\rho_e)_{-1,0}=-g_{\mathrm Q}^*\langle a^\dagger\rangle+O(|g_{\mathrm Q}|^2).
 \label{eq:feqo-quantum-coherence-first-order}
\end{equation}
\end{proposition}
\begin{proof}
吸收末态光向量为 $g_{\mathrm Q}a|\gamma\rangle$，取范数平方给
$|g_{\mathrm Q}|^2\langle a^\dagger a\rangle$；发射对应
$|g_{\mathrm Q}|^2\langle aa^\dagger\rangle$。二阶返回项保证总概率为一。
吸收向量与无交换光向量的重叠给 $g_{\mathrm Q}\langle a\rangle$，发射同理。
混合输入由线性性成立。有限阶余项要求相应光子数矩有限，且显著占据扇区满足
$|g_{\mathrm Q}|\sqrt{n+1}\ll1$；有限失谐则须使用逐边的时间积分。
\end{proof}
平均光子数相同的 Fock、热态与相干态在最低阶人口上可以相同，
但前两者 $\langle a\rangle=0$，相干态保留一阶电子相干。
高阶增益的领先概率为
$P_{+k}=|g_{\mathrm Q}|^{2k}\langle(N)_k\rangle/(k!)^2$。
固定 $\bar n$ 后，热态 $\langle(N)_k\rangle=k!\bar n^k$，相干态为 $\bar n^k$，
因此两种增益概率的比值为 $k!$；同时绝对概率仍受 $(k!)^{-2}$ 抑制。
\begin{figure}[htbp]
 \centering
 \includegraphics[width=0.76\linewidth]{photon_statistics_gain.pdf}
 \caption{相同 $\bar n=3$ 下的高阶增益系数 $P_{+k}/|g_{\mathrm Q}|^{2k}$。
 相干态为 $\bar n^k/(k!)^2$，单模热态为 $\bar n^k/k!$；两者比值为 $k!$。
 图示领先微扰系数，不能把任意大的 $k$ 系数直接当成有限耦合概率。}
 \label{fig:photon-statistics-gain}
\end{figure}


下降阶乘矩来自净增益 $r_{\rm ph}$ 格所要求的 $r_{\rm ph}$ 次光子湮灭；弱耦合只用于截断更高阶的成对往返交换。

电子净增益满足 $\ell\le n$，因为净吸收 $\ell$ 个光子要求 $\hat a^\ell|n\rangle\neq0$；净损失对应产生光子，没有相同的上界。

\subsection{相干光照下，电子何时仍可视为纯态}
\label{sec:electron-quantum-exchange-coherence}

经典 PINEM 给出一个纯电子态。若把同一束相干光完全量子化，电子输出是否仍等于这个纯态，需要对光场求偏迹才能判断。相干态满足 $\hat a|\alpha_{\rm coh}\rangle=\alpha_{\rm coh}|\alpha_{\rm coh}\rangle$，吸收项因此产生熟悉的经典振幅；发射项含有 $\hat a^\dagger$，还会产生不能用这个数代替的分量。下面保留两项，计算电子的密度矩阵、能谱和纯度。光子统计在电子谱中的表现由同一联合态描述。\cite{Dahan2021Science}

在已建立的单模、双边平移、低反冲交换模型中，去掉不影响观测的整体 Magnus 相位，令
\begin{equation*}
 S_Q=\exp(g_{\mathrm Q}\hat b^\dagger\hat a-g_{\mathrm Q}^*\hat b\hat a^\dagger),
 \qquad r_Q=|g_{\mathrm Q}|^2,\qquad \beta=\alpha_{\rm coh}g_{\mathrm Q}.
\end{equation*}
$r_Q$ 是无量纲单电子真空交换强度；它不同于由宏观场增强的 $|\beta|^2$。对电子中心通道 $|0\rangle_e$ 和相干光输入，定义经典 PINEM 纯态
\begin{equation*}
 |\psi_\beta\rangle=S_{\rm tr}(\beta)|0\rangle
 =\sum_\ell J_\ell(2|\beta|)\ee^{\ii\ell\arg\beta}|\ell\rangle.
\end{equation*}
精确约化电子态可写成一组向损失方向平移的纯态混合：
\begin{equation}
 \rho_e^{\rm out}=\sum_{m=0}^{\infty}w_m\,
 \hat b^m|\psi_\beta\rangle\langle\psi_\beta|(\hat b^\dagger)^m,
 \qquad w_m=\ee^{-r_Q}\frac{r_Q^m}{m!}.
 \label{eq:electron-coherent-light-reduced-state}
\end{equation}
每个 $\hat b^m|\psi_\beta\rangle$ 都是把经典输出向能量损失方向移动 $m$ 阶后的归一态，$w_m$ 是它的概率权重。局部推导中会把光的平均相干幅度移去，余下的不同光子数对应这些不同的电子平移。只观测电子时，对正交的光子末态求和，便得到上式的混合。

\begin{proposition}[相干光输入的真空记录分解]
对理想双边交换模型，电子单格与光相干态因子化输入的约化电子态为\eqnrefs{eq:electron-coherent-light-reduced-state}。边带随机变量可写成独立 $L_\beta-M$，其中 $\Pr(L_\beta=j)=J_j^2(2|\beta|)$，$M\sim\operatorname{Poisson}(|g_{\mathrm Q}|^2)$；故均值、方差为\eqnrefs{eq:electron-coherent-light-spectrum-moments}。
\end{proposition}
\begin{proof}
用位移算符 $D(\alpha)=\exp(\alpha\hat a^\dagger-\alpha^*\hat a)$，有 $|\alpha\rangle=D(\alpha)|0\rangle$ 与 $D^\dagger\hat aD=\hat a+\alpha$。因此
\begin{equation*}
 D^\dagger S_QD
 =\exp[g_{\mathrm Q}\hat b^\dagger\hat a-g_{\mathrm Q}^*\hat b\hat a^\dagger
       +\beta\hat b^\dagger-\beta^*\hat b]
 =S_Q S_{\rm tr}(\beta).
\end{equation*}
最后一步使用双边平移的 $[\hat b,\hat b^\dagger]=0$，使量子部分与经典部分对易。对光真空作用时，令 $A=-g_{\mathrm Q}^*\hat b\hat a^\dagger$、$B=g_{\mathrm Q}\hat b^\dagger\hat a$，则 $[A,B]=r_Q I$，正规排序给
\begin{equation*}
 S_Q|0,0\rangle=\ee^{-r_Q/2}\sum_{m=0}^{\infty}
 \frac{(-g_{\mathrm Q}^*)^m}{\sqrt{m!}}|-m,m\rangle.
\end{equation*}
移回的 $D(\alpha)$ 只作用于被取迹的光场，不改变电子约化态；不同 $m$ 的光记录正交，因而得到\eqnrefs{eq:electron-coherent-light-reduced-state}。谱中随机整数可写作 $\ell=L_\beta-M$，其中 $P(L_\beta=j)=J_j^2(2|\beta|)$、$M$ 服从均值 $r_Q$ 的 Poisson 分布且二者独立。利用 $\langle L_\beta\rangle=0$、$\operatorname{Var}L_\beta=2|\beta|^2$，就得到\eqnrefs{eq:electron-coherent-light-spectrum-moments}。
\end{proof}


能谱、平均能量交换和方差随即成为
\begin{equation}
 \begin{aligned}
 P_\ell&=\sum_{m=0}^{\infty}\ee^{-r_Q}\frac{r_Q^m}{m!}
 J_{\ell+m}^2(2|\beta|),\\
 \langle\ell\rangle&=-r_Q,\qquad
 \operatorname{Var}\ell=2|\beta|^2+r_Q.
 \end{aligned}
 \label{eq:electron-coherent-light-spectrum-moments}
\end{equation}
负的均值表示电子向光场净提供能量 $r_Q\hbar\omega$。经典场贡献给出对称展宽，真空项给出额外的损失偏置；这两项在一个概率卷积中同时出现。

为了量化混合程度，使用\emph{纯度} $\operatorname{Tr}\rho^2$。若归一化密度矩阵的本征值为 $w_m$，纯度就是 $\sum_mw_m^2$；只有一个本征值等于一时为纯态，几个非零权重共同存在时纯度小于一。这里的平移态彼此正交：$S_{\rm tr}$ 与 $\hat b$ 对易，而 $\langle0|\hat b^{m-n}|0\rangle=\delta_{mn}$。因此已经找到了本征值 $w_m$，无需另求矩阵对角化，得到
\begin{equation}
 \operatorname{Tr}[(\rho_e^{\rm out})^2]
 =\sum_{m=0}^{\infty}w_m^2
 =\ee^{-2r_Q}I_0(2r_Q),
 \label{eq:electron-coherent-light-purity}
\end{equation}
其中 $I_0(x)=\sum_{m\ge0}(x^2/4)^m/(m!)^2$ 为零阶第一类修正 Bessel 函数。它与振荡型的 $J_0$ 不同。这个纯度在本模型中只依赖 $r_Q$；增大相干光幅度会增加经典调制，却不会自动消除电子留下的真空量子记录。

\begin{exbox}[经典纯态近似的误差怎样随耦合缩小]
以 $|\psi_\beta\rangle$ 为经典参考态，求量子电子态偏离它的大小，并比较“能谱接近”与“量子态接近”两种判据。

前面已证明各平移态彼此正交；未平移项的本征值为 $w_0=\ee^{-r_Q}$。因此与参考纯态的重叠概率为
\begin{equation*}
 F_\beta\equiv\langle\psi_\beta|\rho_e^{\rm out}|\psi_\beta\rangle
 =\ee^{-r_Q}.
\end{equation*}
采用\emph{迹距离} $\mathcal D_{\rm tr}(\rho,\sigma)=\tfrac12\|\rho-\sigma\|_1$，其中迹范数是奇异值之和；对厄米差矩阵，就是本征值绝对值之和。这里差矩阵的本征值是 $w_0-1,w_1,w_2,\ldots$，故无需做无限维对角化：
\begin{equation*}
 \mathcal D_{\rm tr}
 =\frac12\bigl[(1-w_0)+\sum_{m\ge1}w_m\bigr]
 =1-\ee^{-r_Q}.
\end{equation*}
迹距离控制任意一个测量事件的概率偏差。要求该误差不超过 $0<\varepsilon<1$，等价于
\begin{equation*}
 r_Q\le-\ln(1-\varepsilon).
\end{equation*}
小耦合时，$\mathcal D_{\rm tr}=r_Q+O(r_Q^2)$，纯度则为
$\operatorname{Tr}[(\rho_e^{\rm out})^2]=1-2r_Q+3r_Q^2+O(r_Q^3)$。

相反，量子方差相对于经典方差的增量在 $\beta\ne0$ 时为
\begin{equation*}
 \frac{\operatorname{Var}\ell-2|\beta|^2}{2|\beta|^2}
 =\frac{r_Q}{2|\beta|^2}
 =\frac1{2|\alpha_{\rm coh}|^2}.
\end{equation*}
增大平均光子数会让这个\emph{相对方差误差}缩小，却不改变固定 $r_Q$ 下的迹距离。要同时恢复完整的经典电子纯态，应取 $g_{\mathrm Q}\to0$、$|\alpha_{\rm coh}|\to\infty$ 且 $\beta$ 固定。在这个极限中，经典调制保持不变，而真空交换权重趋于零。
\end{exbox}

弱 Fock 交换还需要单独检查玻色增强。对确定光子数 $n$，截断单量子概率需要 $r_Q(n+1)\ll1$；对有宽尾的光子分布，应检查相应高阶光子矩及实际占据尾。于是 $P_{+1}\simeq r_Q\langle n\rangle$ 与高阶阶乘矩公式都不是只由 $r_Q\ll1$ 保证。经典场和量子场两种近似条件应分别报告。


\subsection{联合振幅怎样决定纠缠与条件光态}
\label{chap:feqo-conditioning}
求得联合振幅以后，电子谱仍只保存其中的一部分信息。
若不同电子通道留下不同光学记录，对光求迹会减小电子相干；
若再测量这些记录，可以选择条件电子态。
这两个操作都由同一系数矩阵计算。

\begin{definition}[联合纯态的 Schmidt 权重与纠缠熵]
对可分 Hilbert 空间中的归一化联合纯态
$|\Psi\rangle=\sum_{\ell,n}C_{\ell n}|\ell,n\rangle$，
令 $\rho_e=\operatorname{Tr}_{\rm ph}|\Psi\rangle\langle\Psi|$。
其非零本征值 $\lambda_k$ 称为 Schmidt 权重，$\sqrt{\lambda_k}$ 是 Schmidt 系数。
纠缠熵定义为
\begin{equation}
S_e=-\operatorname{Tr}(\rho_e\ln\rho_e)
=-\sum_k\lambda_k\ln\lambda_k,\qquad0\ln0=0.
\label{eq:feqo-entanglement-entropy}
\end{equation}
允许 $S_e=+\infty$；这仍是熵的定义，不需以任意截断值代替它。
\end{definition}

\begin{theorem}[联合振幅的奇异值与可分性]
系数矩阵 $C$ 是 Hilbert--Schmidt 算符，且存在展开
\begin{equation}
|\Psi\rangle=\sum_k\sqrt{\lambda_k}|u_k\rangle_e|v_k\rangle_{\rm ph},
\qquad\sum_k\lambda_k=1.
\label{eq:feqo-schmidt-decomposition}
\end{equation}
联合纯态可分当且仅当 $C$ 秩为一，亦即所有非零列共线。
若有限系统的较小维数为 $d$，则 $0\le S_e\le\ln d$，
上界当且仅当 $d$ 个权重全为 $1/d$。
\end{theorem}
\begin{proof}
归一化给 $\sum_{\ell,n}|C_{\ell n}|^2=1$，因此 $C$ 为紧的 Hilbert--Schmidt 算符。
其奇异值分解 $C=\sum_k s_k u_k w_k^\dagger$ 满足 $\sum_ks_k^2=1$。
把矩阵元代回联合态，光侧向量取 $|v_k\rangle=\sum_n(w_k)_n^*|n\rangle$，
即可得到 $\lambda_k=s_k^2$ 的展开。
对任一子系统求迹，另一侧正交性给相同的非零本征值。
只有一个奇异值时是积态；反向把积态系数写为 $C_{\ell n}=u_\ell v_n$ 就证明秩一。
\begin{equation}
\operatorname{Tr}\rho_e^2
=\sum_{n,m}|\langle c_n|c_m\rangle|^2\le
\sum_{n,m}\|c_n\|^2\|c_m\|^2=1,
\label{eq:feqo-entanglement-column-test}
\end{equation}
其中 $c_n$ 是第 $n$ 列；等号要求所有非零列满足 Cauchy--Schwarz 等号条件，即共线。
有限维熵上界由
$\sum_k\lambda_k\ln(d\lambda_k)\ge0$ 得到，等号要求均匀权重。
对于混合联合态，局部混合可以来自经典概率；局部熵本身不能直接称为纠缠熵。
\end{proof}

\begin{exbox}[真空输入：Poisson 交换与纠缠熵的完整计算]
取电子单格、光真空的因子化输入，求联合末态、电子纯度与纠缠熵，并判断相同的 Poisson 数分布是否足以辨认光相干态。

\textbf{解。} 令 $\zeta=-g_{\rm Q}^*$、$\mu=|g_{\rm Q}|^2$。
由位移算符的正规排序作用于真空，单模输出为
\begin{equation}
|\Psi\rangle=\ee^{-\mu/2}\sum_{m=0}^\infty
\frac{\zeta^m}{\sqrt{m!}}|-m,m\rangle.
\label{eq:feqo-vacuum-joint-state}
\end{equation}
负标号只属于电子能量格，光子数 $m$ 始终非负。
两侧基分别正交，所以
\begin{align}
\lambda_m&=\ee^{-\mu}\mu^m/m!,\label{eq:feqo-vacuum-schmidt}\\
\rho_e&=\sum_m\lambda_m|-m\rangle\langle-m|,
\label{eq:feqo-vacuum-reduced-electron}\\
\operatorname{Tr}\rho_e^2&=\ee^{-2\mu}I_0(2\mu),
\label{eq:feqo-vacuum-purity}\\
S_e&=\mu(1-\ln\mu)+\sum_m\lambda_m\ln(m!).
\label{eq:feqo-vacuum-entropy}
\end{align}
第一式的概率和是 $\ee^{-\mu}\sum_m\mu^m/m!=1$。当 $\mu=0$ 时，熵式中的 $\mu\ln\mu$ 按连续极限取零，输出仍为积态。
纯度式由 $I_0(2\mu)=\sum_m\mu^{2m}/(m!)^2$ 得到。
自然对数下，$\mu=0.5,1,4$ 的熵分别为
$0.92764,1.30484,2.08667$；$\mu=1$ 时纯度为 $0.30851$。
大 $\mu$ 时主峰宽度为 $O(\sqrt\mu)$，局部高斯渐近给
$S_e=\tfrac12\ln(2\pi\ee\mu)+O(\mu^{-1})$。
这表示对数增长，并没有在某个有限值饱和。

若以电子损失 $m\hbar\omega$ 为条件，光场精确成为 $|m\rangle$。
无条件光态的数分布虽为 Poisson，密度矩阵却在数态基对角，故
$\langle a\rangle=0$；相干态 $|\alpha\rangle$ 可具有相同数分布却有
$\langle a\rangle=\alpha$。数分布相同并不等于量子态相同。
\end{exbox}

\begin{exbox}[光学记录重叠、纠缠和量子擦除]
对归一化态 $|\Psi\rangle=\sqrt p|0\rangle|u\rangle
+\ee^{\ii\phi}\sqrt{1-p}|1\rangle|v\rangle$，令 $s=\langle v|u\rangle$。
设 $0\le p\le1$，$|u\rangle,|v\rangle$ 为归一化光学记录。求电子相干与纠缠怎样依赖 $s$，再判断读取光学记录能否恢复条件电子相干。

\textbf{解。} 对光求偏迹后，电子的非对角元为
\begin{equation}
(\rho_e)_{01}=\sqrt{p(1-p)}\ee^{-\ii\phi}s,
\quad\operatorname{Tr}\rho_e^2=p^2+(1-p)^2+2p(1-p)|s|^2.
\label{eq:feqo-which-path-overlap}
\end{equation}
约化态两个本征值为
$\lambda_\pm=[1\pm\sqrt{1-4p(1-p)(1-|s|^2)}]/2$，
所以熵可用这两个权重直接计算。
光学记录相同给纯电子态；记录正交且 $p=1/2$ 给熵 $\ln2$。

若反冲与选择规则已受控地隔离联合二态
$|0,1\rangle,|1,0\rangle$，一次半转移可制备
$(|0,1\rangle-\ii|1,0\rangle)/\sqrt2$。
这是一个受限模型，不能在无限无反冲交换链上任意扔掉 $|-1,2\rangle$ 来获得。
读光到 $|\pm\rangle=(|0\rangle\pm|1\rangle)/\sqrt2$，每个结果概率为 $1/2$，
电子条件态分别为 $(|0\rangle\mp\ii|1\rangle)/\sqrt2$。
\begin{equation}
\widetilde\rho_{e|y}=\operatorname{Tr}_{\rm ph}
[(I\otimes P_y)\rho_{e,\rm ph}(I\otimes P_y)],
\quad p_y=\operatorname{Tr}\widetilde\rho_{e|y},
\quad\rho_{e|y}=\widetilde\rho_{e|y}/p_y.
\label{eq:feqo-conditional-electron-state}
\end{equation}
按两个概率加权，条件纯态的相干互相抵消，恢复无条件混态。
恢复条件干涉需要保留结果标签，而不是对两种归一化态不加权求和。
\end{exbox}

\begin{exbox}[依次通过两个腔：总损失相同，光学记录可以不同]
保留两个互相独立、同频的模式 $a_1,a_2$，输入都是真空。
电子沿用同一个双边格，两个脉冲的交换参数为 $g_1,g_2$。
假定脉冲间的反冲色散相位可忽略，令
$X_j=g_jb^\dagger a_j-g_j^*ba_j^\dagger$。
求电子总损失的分布，以及给定总损失后两个腔的联合光态；取 $\mu=|g_1|^2+|g_2|^2>0$，使非真空结果有可能发生。

\textbf{解。} 不同模式算符对易，且 $bb^\dagger=b^\dagger b=I$，故 $[X_1,X_2]=0$。
两个相互作用可以先分别正规排序，再作用于真空，得到
\begin{equation}
|\Psi\rangle=\ee^{-\mu/2}\sum_{m_1,m_2\ge0}
\frac{\zeta_1^{m_1}\zeta_2^{m_2}}{\sqrt{m_1!m_2!}}
|-(m_1+m_2),m_1,m_2\rangle,
\quad\zeta_j=-g_j^*,\quad\mu=|g_1|^2+|g_2|^2.
\label{eq:feqo-two-cavity-vacuum-joint}
\end{equation}
两个光子数的无条件联合计数概率是两个独立 Poisson 分布；
电子却只记录总损失 $n=m_1+m_2$，其概率为 $\ee^{-\mu}\mu^n/n!$。
\begin{equation}
|\psi_{12|n}\rangle=\sum_{m=0}^n
\sqrt{\binom nm}\left(\frac{\zeta_1}{\sqrt\mu}\right)^m
\left(\frac{\zeta_2}{\sqrt\mu}\right)^{n-m}|m,n-m\rangle.
\label{eq:feqo-two-cavity-heralded-binomial}
\end{equation}
这是总光子数为 $n$ 的纯态，数分布为二项分布。
归一化由二项定理给出；相对复相位来自两个腔的耦合相位。
$n=1$ 且两腔耦合模长相同时，电子后选择制备
$(\ee^{\ii\arg\zeta_1}|1,0\rangle+
\ee^{\ii\arg\zeta_2}|0,1\rangle)/\sqrt2$，
在两个腔之间有一个量子比特的纠缠，尽管电子没有分辨光子来自哪一个腔。
若两腔间色散不可忽略，需在两个交换算符中间插入传播子；
不同交换路径将获得不同相位，不能继续使用上述可交换指数。
\end{exbox}

\begin{proposition}[条件光子计数与二阶关联]
设 $0<\eta=|g_1|^2/\mu<1$。
在上式条件态中，$\langle N_1\rangle=n\eta$、
$\langle N_2\rangle=n(1-\eta)$，且
\begin{equation}
\langle N_1N_2\rangle=n(n-1)\eta(1-\eta),\qquad
g_{12}^{(2)}=\frac{\langle N_1N_2\rangle}
{\langle N_1\rangle\langle N_2\rangle}=1-\frac1n\quad(n\ge1).
\label{eq:feqo-heralded-hbt-correlation}
\end{equation}
若不对电子后选择，两个模式的计数独立，$g_{12}^{(2)}=1$。
\end{proposition}
\begin{proof}
对二项变量 $M$，$\mathbb EM=n\eta$、$\mathbb E[M(M-1)]=n(n-1)\eta^2$。
由于 $N_2=n-M$，
$\mathbb E[M(n-M)]=n\mathbb EM-\mathbb EM^2
=n(n-1)\eta(1-\eta)$，得到归一化关联。
当 $n=0$ 时两个分母为零，归一化 $g^{(2)}$ 没有定义，不能写成零。
对单模数态 $|n\rangle$，同样由 $(a^\dagger)^2a^2=N(N-1)$ 得
$g^{(2)}(0)=1-1/n$（$n\ge1$）。
\end{proof}

\begin{exbox}[光子数尾部怎样控制电子计数核]
对一般数分布 $p_n$，已建立的前向核满足
$P_\ell=\sum_np_n\mathsf K_{\ell n}$、$0\le\mathsf K_{\ell n}\le1$。
求截断 $n>N$ 后的计数误差，并比较数态对角输入与相干输入在联合态层面的误差。

\textbf{解。} 保留 $n\le N$ 而不重新归一化，遗漏概率为 $\eta_N=\sum_{n>N}p_n$，于是
$0\le P_\ell-P_\ell^{(N)}\le\eta_N$。
若输入是数态对角混合，对截断输入重新归一化后，原输入与该输入的迹距离就是 $\eta_N$；
CPTP 演化和测量不会增大这个距离，因此任意完整计数分布的总变差也至多为 $\eta_N$。
\[
\eta_N^{\rm th}=\left(\frac{\bar n}{1+\bar n}\right)^{N+1},\qquad
\eta_N^{\rm coh}=1-\ee^{-\bar n}\sum_{n=0}^{N}\frac{\bar n^n}{n!}.
\]
例如平均数为一的热态，要求 $\eta_N<10^{-6}$ 时取 $N=19$ 即可。
对于相干输入，保留数态振幅并重新归一化得到纯态 $|\alpha_N\rangle$，
其重叠满足 $|\langle\alpha|\alpha_N\rangle|^2=1-\eta_N$，所以纯态迹距离是 $\sqrt{\eta_N}$。
单能电子的能谱只依赖输入数分布，仍可使用 $\eta_N$ 的人口核误差；
相位敏感的联合测量或条件态一般只能先使用 $\sqrt{\eta_N}$ 的态误差界。
这个界来自输入尾概率，不需要把未计算的每条高能响应核另作经验估计。
\end{exbox}

\subsection{热光输入：由生成函数求完整交换分布}
\label{sec:electron-thermal-exchange-analytic}

相干光的平均光子数确定以后，仍有许多不同光态具有同样的平均能量。热光就是其中一种。它没有固定光学相位，光子数也遵循与相干态不同的分布。为了求这些差别怎样传给电子，保持上一节的单模低反冲交换算符 $S_Q$ 不变，仍取电子初态 $|0\rangle_e$，只把光场输入换为热态
\begin{equation*}
 \rho_{\rm th}=\sum_{n=0}^\infty
 \frac{\bar n_{\rm th}^{\,n}}{(1+\bar n_{\rm th})^{n+1}}|n\rangle\langle n|,
 \qquad \bar n_{\rm th}=\frac1{\exp[\hbar\omega/(k_{\rm B}T_{\rm ph})]-1}.
\end{equation*}
$T_{\rm ph}$ 是该模式的温度，$\bar n_{\rm th}$ 是无量纲平均光子数。热态在光子数基中为对角混合，没有确定的光学参考相位。热光与相干光对自由电子谱的不同作用可由同一 QPINEM 模型比较\cite{Dahan2021Science}。

出射电子可以占据许多边带。求这组概率时，可以利用 Fourier 级数把所有 $P_\ell$ 合成一个函数，称为概率分布的\emph{特征函数}：
\begin{equation*}
 \chi_e(\vartheta)=\sum_{\ell\in\mathbb Z}P_\ell\ee^{\ii\ell\vartheta}
 =\operatorname{Tr}[\rho_e^{\rm out}\ee^{\ii\vartheta\hat\Lambda_e}],
 \qquad \hat\Lambda_e|\ell\rangle=\ell|\ell\rangle.
\end{equation*}
其中 $\hat\Lambda_e$ 是阶数算符，$\vartheta$ 是无量纲 Fourier 变量。它只是用来整理概率，不代表额外的电子动力学。因为 $\partial_\vartheta\ee^{\ii\ell\vartheta}=\ii\ell\ee^{\ii\ell\vartheta}$，在原点求导便有
\[
\langle\ell\rangle=-\ii\chi_e'(0),\qquad
\langle\ell^2\rangle=-\chi_e''(0),\qquad \chi_e(0)=1.
\]
对数的导数则自动去掉低阶矩的乘积：
\[
-\ii(\ln\chi_e)'(0)=\langle\ell\rangle,\qquad
-(\ln\chi_e)''(0)=\langle\ell^2\rangle-\langle\ell\rangle^2.
\]
这些对数导数称为\emph{累积量}，前两阶就是均值和方差。用本节末尾的算符代数求热态平均，得到
\begin{equation}
 \chi_e(\vartheta)=
 \exp\!\left\{\lambda_+(\ee^{\ii\vartheta}-1)
             +\lambda_-(\ee^{-\ii\vartheta}-1)\right\},
 \qquad
 \lambda_+=r_Q\bar n_{\rm th},\quad
 \lambda_-=r_Q(1+\bar n_{\rm th}).
 \label{eq:electron-thermal-characteristic}
\end{equation}
该结果对本交换模型中任意有限 $r_Q=|g_{\mathrm Q}|^2$ 成立，不是只保留单量子的微扰式。$\lambda_\pm$ 是无量纲统计参数；形式上等同于两个独立 Poisson 变量 $N_\pm$ 的均值，净电子阶数为 $\ell=N_+-N_-$。这是一种\emph{末态分布的等价表示}，不自动意味着真实连续时间动力学是两个无记忆跳跃过程。

\begin{proposition}[热光输入的交换特征函数]
对电子单格与有限平均光子数 $\bar n_{\rm th}$ 的热态因子化输入，理想双边交换给\eqnrefs{eq:electron-thermal-characteristic}。其阶数分布等价于两个独立 Poisson 变量之差，吸收、发射均值分别为 $r_Q\bar n_{\rm th}$ 与 $r_Q(\bar n_{\rm th}+1)$；约化态在电子格点基中对角。
\end{proposition}
\begin{proof}
定义 $B=\hat b^\dagger\hat a$。双边平移给 $[B,B^\dagger]=I$，因此可定义位移形式
$D_B(\zeta)=\exp(\zeta B^\dagger-\zeta^*B)$，并写 $S_Q=D_B(-g_{\mathrm Q}^*)$。
电子计数旋转满足
\begin{equation*}
 \ee^{\ii\vartheta\hat\Lambda_e}S_Q\ee^{-\ii\vartheta\hat\Lambda_e}
 =D_B(-g_{\mathrm Q}^*\ee^{-\ii\vartheta}).
\end{equation*}
由于初态 $|0\rangle_e$ 的阶数为零，$\chi_e$ 可改写为初态上 $S_Q^\dagger S_Q(\vartheta)$ 的期望。使用
\begin{equation*}
 D_B(\alpha)D_B(\beta)
 =\ee^{(\alpha\beta^*-\alpha^*\beta)/2}D_B(\alpha+\beta)
\end{equation*}
得到
\begin{equation*}
 \chi_e(\vartheta)=\ee^{-\ii r_Q\sin\vartheta}
 \left\langle D_B[-g_{\mathrm Q}^*(\ee^{-\ii\vartheta}-1)]\right\rangle_{0,\rm th}.
\end{equation*}
接下来显式求平均。正规排序给
$D_B(\zeta)=\ee^{-|\zeta|^2/2}\ee^{\zeta B^\dagger}\ee^{-\zeta^*B}$。
在 $|0,n\rangle$ 上，只有同次数的 $B^\dagger$ 与 $B$ 能返回原电子阶数，因此
\begin{equation*}
 \langle0,n|D_B(\zeta)|0,n\rangle
 =\ee^{-|\zeta|^2/2}\sum_{k=0}^n
 \frac{(-|\zeta|^2)^k}{(k!)^2}\frac{n!}{(n-k)!}
 =\ee^{-|\zeta|^2/2}L_n(|\zeta|^2).
\end{equation*}
$L_n$ 是普通 Laguerre 多项式，上一式的有限和也可作为其定义。令 $q_{\rm th}=\bar n_{\rm th}/(1+\bar n_{\rm th})$，用几何级数加权后的恒等式
\begin{equation*}
 (1-q_{\rm th})\sum_{n=0}^\infty q_{\rm th}^nL_n(x)
 =\exp[-\bar n_{\rm th}x]
\end{equation*}
这个恒等式可由几何级数逐项求导验证：对 $\sum_{n\ge0}q^n=(1-q)^{-1}$ 求 $k$ 阶导数并乘 $q^k$，有
\begin{equation*}
 (1-q)\sum_{n=k}^{\infty}q^n\frac{n!}{(n-k)!}
 =k!\left(\frac{q}{1-q}\right)^k.
\end{equation*}
将它代入 $L_n$ 的有限和，剩下 $\sum_k[-xq/(1-q)]^k/k!$，正是指数级数。于是得到 $\langle D_B(\zeta)\rangle_{0,\rm th}
=\exp[-(\bar n_{\rm th}+1/2)|\zeta|^2]$。
最后代入 $|\zeta|^2=2r_Q(1-\cos\vartheta)$，
\begin{equation*}
 \ln\chi_e=-\ii r_Q\sin\vartheta
 -r_Q(2\bar n_{\rm th}+1)(1-\cos\vartheta),
\end{equation*}
按 $\ee^{\pm\ii\vartheta}$ 重组即得\eqnrefs{eq:electron-thermal-characteristic}。热态对光学相位旋转不变，初始电子又在单一阶数上；交换数守恒遂使电子约化态在阶数基中保持对角。这里的 $P_\ell$ 因而给出了完整约化电子态，而非只给其一部分矩阵元。
\end{proof}


直接对 $\ln\chi_e$ 求导得
\begin{equation}
 \langle\ell\rangle=\lambda_+-\lambda_-=-r_Q,\qquad
 \operatorname{Var}\ell=\lambda_++\lambda_-=r_Q(2\bar n_{\rm th}+1).
 \label{eq:electron-thermal-moments}
\end{equation}
热光增强增益和损失两侧的涨落，但净损失偏置仍来自真空项。取 $\bar n_{\rm th}=|\alpha_{\rm coh}|^2$，热光与前节相干光的均值和方差完全相同；要区分两者，必须读更高统计矩、谱形或相干，不能只读展宽。

\begin{exbox}[解析恢复热交换概率并比较四阶累积量]
利用\eqnrefs{eq:electron-thermal-characteristic}求各阶 $P_\ell$、正温增益--损失比和四阶累积量，并与同平均光子数的相干态比较。

\textbf{解。} 对 $\ell\ge0$，条件 $N_+-N_-=\ell$ 等价于 $N_+=m+\ell,N_-=m$，故
\begin{equation*}
 P_\ell=\ee^{-(\lambda_++\lambda_-)}
 \sum_{m=0}^\infty\frac{\lambda_+^{m+\ell}\lambda_-^m}{(m+\ell)!m!}.
\end{equation*}
用第一类修正 Bessel 函数的定义
$I_j(2x)=\sum_{m\ge0}x^{2m+j}/[m!(m+j)!]$，并对负 $\ell$ 交换两参数，得到
\begin{equation*}
 P_\ell=\ee^{-(\lambda_++\lambda_-)}
 \left(\frac{\lambda_+}{\lambda_-}\right)^{\ell/2}
 I_{|\ell|}(2\sqrt{\lambda_+\lambda_-}),\qquad \bar n_{\rm th}>0.
\end{equation*}
它称为两个 Poisson 变量之差的 Skellam 分布。零温不能直接把零代入负幂；应回到求和式或特征函数，得到
$P_{-m}=\ee^{-r_Q}r_Q^m/m!$（$m\ge0$），其余增益概率为零。对 $r_Q>0$、任意正温度，还可直接相除：
\begin{equation*}
 \frac{P_{+\ell}}{P_{-\ell}}
 =\left(\frac{\bar n_{\rm th}}{1+\bar n_{\rm th}}\right)^\ell
 =\ee^{-\ell\hbar\omega/(k_{\rm B}T_{\rm ph})},\qquad\ell>0.
\end{equation*}
这给出了不依赖总体耦合强度的增益--损失比。

再定义四阶累积量 $\mathcal C_4=\langle(\ell-\langle\ell\rangle)^4\rangle-3(\operatorname{Var}\ell)^2$。它是对高斯统计之外结构的一个度量。由\eqnrefs{eq:electron-thermal-characteristic}，
\begin{equation*}
 \mathcal C_4^{\rm(th)}=r_Q(2\bar n_{\rm th}+1).
\end{equation*}
相干光的经典 Bessel 部分具有特征函数 $J_0[4|\beta|\sin(\vartheta/2)]$，真空 Poisson 损失与它独立。将 $\ln J_0$ 展开至 $\vartheta^4$，得到
\begin{equation*}
 \mathcal C_4^{\rm(coh)}=r_Q+2|\beta|^2-6|\beta|^4.
\end{equation*}
在相同平均光子数、$|\beta|^2=r_Q\bar n_{\rm th}$ 下，两种方差相同，四阶累积量却相差 $6r_Q^2\bar n_{\rm th}^2$。这说明为什么只拟合一个高斯宽度会丢失光子统计信息。
\end{exbox}



\subsection{有限反冲下的联合量子格}\label{subsec:qpinem-finite-recoil-lattice-dp46}

\eqnrefs{eq:qpinem_S} 的位移算符闭式还使用了电子能量梯的平移对称性。若电子色散曲率已经使相邻交换边可分辨，则不能把反冲只作为输出相位补在散射之后，而必须把它放回联合电子--光子动力学。对初始光子数为 $n$ 的守恒扇区，取基底
\begin{equation*}
 |\ell;n\rangle\equiv|\ell\rangle_e\otimes|n-\ell\rangle_{\rm ph},
 \qquad \ell\le n,
\end{equation*}
并定义相对于理想均匀能量梯的通道失谐
\begin{equation}
\boxed{
 \delta_\ell
 \equiv\frac{\mathcal E_\ell-E_0-\ell\hbar\omega}{\hbar}.}
\label{eq:qpinem_recoil_channel_detuning_dp46}
\end{equation}
这里 $\mathcal E_\ell=E(p_0+\ell\hbar \kappa_z)$ 始终取真实电子色散；在二阶窄带展开下，若中心同步条件 $v_0\kappa_z=\omega$ 成立，则
\begin{equation}
 \delta_\ell=\ell^2\omega_{\rm r}+O(\ell^3),
 \qquad
 \delta_{\ell+1}-\delta_\ell=(2\ell+1)\omega_{\rm r}+O(\ell^2).
\label{eq:qpinem_recoil_detuning_quadratic_dp46}
\end{equation}

有限反冲方程采用\eqnrefs{eq:feqo-multimode-coupling-frame}所说明的联合旋转格，而不是已经旋去真实自由能差的全相互作用绘景。定义该格中的交换速率
\begin{equation}
 \upsilon_{\rm rot}(t)\equiv-\ii\mathcal G_0(t)
 =\upsilon(t)\ee^{-\ii\Delta_0(t-t_i)}.
 \label{eq:feqo-single-mode-coupling-frame}
\end{equation}
这里共同失谐及其逐边变化全部由 $\delta_\ell$ 承担。若共同失谐恰好为零，两种速率相同；非同步时同时使用含 $\Delta_0$ 的 $\upsilon$ 与完整 $\delta_\ell$ 会把共同失谐算两次。

把联合态写成
\begin{equation*}
 |\Psi_n(t)\rangle=\sum_{\ell\le n}c_\ell^{(n)}(t)|\ell;n\rangle,
 \qquad
 \sum_{\ell\le n}|c_\ell^{(n)}|^2=1,
\end{equation*}
则单模交换与真实电子色散共同给出离散 Schr\"odinger 方程
\begin{equation}
\boxed{
 \ii\dot c_\ell^{(n)}
 =\delta_\ell c_\ell^{(n)}
 +\ii\upsilon_{\rm rot}(t)\sqrt{n-\ell+1}\,c_{\ell-1}^{(n)}
 -\ii\upsilon_{\rm rot}^*(t)\sqrt{n-\ell}\,c_{\ell+1}^{(n)}.}
\label{eq:qpinem_finite_recoil_coupled_lattice_dp46}
\end{equation}
这是一维非均匀紧束缚链：$\delta_\ell$ 是电子反冲产生的在位能，$\sqrt{n-\ell}$ 是光子 Fock 下界产生的非均匀跃迁矩阵元。二者来源完全不同，因此不能合并成一个经验“有效耦合”。\eqnrefs{eq:qpinem_finite_recoil_coupled_lattice_dp46}也是数值计算有限反冲 QPINEM 的最小闭合模型；给定 $\upsilon_{\rm rot}(t)$、$p_0$、$\kappa_z$ 与初始光子统计后即可直接积分得到联合末态。

\begin{proposition}[交换守恒扇区中的反冲格]
设交换数 $\Lambda_e+N_{\rm ph}=n$ 固定，电子色散为实对角算符。正交基 $|\ell;n\rangle=|\ell\rangle_e|n-\ell\rangle_{\rm ph}$，$\ell\le n$，中的精确振幅方程为\eqnrefs{eq:qpinem_finite_recoil_coupled_lattice_dp46}。二次同步色散的在位失谐为\eqnrefs{eq:qpinem_recoil_detuning_quadratic_dp46}，更高阶余项仍须保留。
\end{proposition}
\begin{proof}
先把“固定交换量扇区”写成真正的张量积基底。正文定义的守恒算符为
\begin{equation*}
\mathcal N_{\rm ex}=\Lambda_e+N_{\rm ph},
\end{equation*}
其中 \(\Lambda_e|\ell\rangle_e=\ell|\ell\rangle_e\)，
\(N_{\rm ph}|m\rangle_{\rm ph}=m|m\rangle_{\rm ph}\)。固定本征值 \(n\) 后定义
\begin{equation*}
|\ell;n\rangle
\equiv
|\ell\rangle_e\otimes|n-\ell\rangle_{\rm ph},
\qquad n-\ell\ge0.
\end{equation*}
于是
\begin{equation*}
\mathcal N_{\rm ex}|\ell;n\rangle
=[\ell+(n-\ell)]|\ell;n\rangle
=n|\ell;n\rangle.
\end{equation*}
这个定义把电子“上升一格”和光子“减少一个”写进了同一个基矢，而不是事后再猜平方根因子。

在相对于线性同步梯 \(E_0+\ell\hbar\omega\) 的旋转表象中，真实电子色散留下
\begin{equation*}
H_{\rm rec}
=\hbar\sum_\ell\delta_\ell|\ell\rangle_e\langle\ell|,
\end{equation*}
交换项为
\begin{equation*}
H_Q(t)
=\ii\hbar\left[
\upsilon_{\rm rot}(t)\hat b^\dagger\hat a
-\upsilon_{\rm rot}^*(t)\hat b\hat a^\dagger
\right].
\end{equation*}
现在逐项作用。对左邻通道，
\begin{align*}
\hat b^\dagger\hat a|\ell-1;n\rangle
&=\hat b^\dagger\hat a
\left(|\ell-1\rangle_e\otimes|n-\ell+1\rangle_{\rm ph}\right)\\
&=\sqrt{n-\ell+1}\,
|\ell\rangle_e\otimes|n-\ell\rangle_{\rm ph}\\
&=\sqrt{n-\ell+1}\,|\ell;n\rangle.
\end{align*}
对右邻通道，
\begin{align*}
\hat b\hat a^\dagger|\ell+1;n\rangle
&=\hat b\hat a^\dagger
\left(|\ell+1\rangle_e\otimes|n-\ell-1\rangle_{\rm ph}\right)\\
&=\sqrt{n-\ell}\,|\ell;n\rangle.
\end{align*}
因此 \(\sqrt{n-\ell+1}\) 与 \(\sqrt{n-\ell}\) 分别是玻色湮灭、产生算符在实际光子数上的矩阵元；它们不是人为放进耦合常数的修正。

展开固定扇区态
\begin{equation*}
|\Psi_n(t)\rangle
=\sum_{\ell\le n}c_\ell^{(n)}(t)|\ell;n\rangle
\end{equation*}
并代入
\(\ii\hbar\partial_t|\Psi_n\rangle=(H_{\rm rec}+H_Q)|\Psi_n\rangle\)。左乘 \(\langle\ell;n|\) 后，正交性逐项给出
\begin{align*}
\ii\dot c_\ell^{(n)}
={}&\delta_\ell c_\ell^{(n)}
+\ii\upsilon_{\rm rot}(t)\sqrt{n-\ell+1}\,c_{\ell-1}^{(n)}\\
&-\ii\upsilon_{\rm rot}^*(t)\sqrt{n-\ell}\,c_{\ell+1}^{(n)},
\end{align*}
即正文\eqnrefs{eq:qpinem_finite_recoil_coupled_lattice_dp46}。

最后把失谐本身接回电子色散。若
\begin{equation*}
\mathcal E_\ell
=E_0+\ell\hbar v_0\kappa_z
+\ell^2\hbar\omega_{\rm r}+R_\ell^{(3)},
\end{equation*}
且中心同步 \(v_0\kappa_z=\omega\)，则
\begin{align*}
\delta_\ell
&\equiv\frac{\mathcal E_\ell-E_0-\ell\hbar\omega}{\hbar}\\
&=\ell^2\omega_{\rm r}+\frac{R_\ell^{(3)}}{\hbar},
\end{align*}
从而得到正文\eqnrefs{eq:qpinem_recoil_detuning_quadratic_dp46}。
\end{proof}


对实际显著占据的电子窗口 $|\ell|\le N_{\rm occ}$，反冲是否能在作用时间 $T$ 内被分辨可用
\begin{equation}
\boxed{
 \epsilon_{\rm Qrec}^{(T)}
 \equiv\max_{|\ell|\le N_{\rm occ}}
 |\delta_{\ell+1}-\delta_\ell|T
 \simeq(2N_{\rm occ}+1)\omega_{\rm r}T.}
\label{eq:qpinem_recoil_time_control_dp46}
\end{equation}
在常幅作用段内，若还要求不同交换边的失谐小于其局域 Rabi 速率，则对给定光子数扇区只在允许吸收的边上定义
\begin{equation}
\boxed{
 \epsilon_{\rm Qrec}^{(\Omega)}(n)
 \equiv
 \max_{\substack{|\ell|\le N_{\rm occ}\\\ell<n}}
 \frac{|\delta_{\ell+1}-\delta_\ell|}
 {|\upsilon_{\rm rot}|\sqrt{n-\ell}}
 \ll1.}
\label{eq:qpinem_recoil_rabi_control_dp46}
\end{equation}
这里不把 $\ell=n$ 的光子真空上边界放入分母，也不在脉冲关闭时使用瞬时速率比值。\eqnrefs{eq:qpinem_recoil_rabi_control_dp46}适合判断持续交换是否能形成近共振振荡；短弱脉冲即使失谐大于耦合率，仍可能因失谐累计相位很小而接近均匀梯传播。实际使用位移闭式，应以输出态的累计误差为准，而不是把两个尺度比较都当作必要条件。

把逐边失谐放回全相互作用绘景后，可直接估计位移模型遗漏的振幅。令 $\tau=t-t_i$，真实第 $\ell$ 条吸收边的速率为
\[
 v_\ell(t)=-\ii\mathcal G_\ell(t)\ee^{\ii\Delta_\ell\tau},\qquad
 \Delta_\ell=\delta_{\ell+1}-\delta_\ell,
\]
均匀模型以 $v_0(t)=\upsilon(t)$ 代替所有 $v_\ell$。这里只比较逐边变化；物理格边界和被舍环境模式的误差需另外加入。设 $P_W$ 投影到选定电子工作窗，$\mathcal E_W$ 包含与该窗中至少一个格点相接的交换边。由 $|\ee^{\ii x}-1|\le|x|$，每条相关边的速率差满足
\begin{equation}
 \max_{\ell\in\mathcal E_W}|v_\ell(t)-v_0(t)|
 \le a_W(t)\equiv
 \max_{\ell\in\mathcal E_W}|\mathcal G_\ell(t)-\mathcal G_0(t)|
 +|\mathcal G_0(t)|\tau
  \max_{\ell\in\mathcal E_W}|\Delta_\ell-\Delta_0|.
 \label{eq:feqo-quantum-edge-rate-error}
\end{equation}
采用绝对速率差，便不会在脉冲零点或共同矩阵元为零时除以零。

\begin{proposition}[逐边变化与玻色数矩的传播误差]
令 $U$ 为真实逐边交换传播子，$U_0$ 为共同速率的双边位移模型传播子，二者存在于同一联合空间中。对归一化入射态，记 $\Psi_0(t)=U_0(t,t_i)\Psi_i$，$R(t)=H(t)-H_0(t)$。若相应 Duhamel 积分与数矩有限，则
\begin{equation}
\begin{aligned}
 \|(U-U_0)\Psi_i\|
 &\le2\int_{t_i}^{t_f}a_W(t)
       \sqrt{\langle N_{\rm ph}\rangle_{0,t}+1}\,\dd t
       +\epsilon_{W,\rm out},\\
 \epsilon_{W,\rm out}
 &\equiv\frac1\hbar\int_{t_i}^{t_f}
       \|R(t)(I-P_W)\Psi_0(t)\|\dd t.
\end{aligned}
\label{eq:feqo-quantum-ladder-propagation-error}
\end{equation}
这里数矩属于参考演化的实际光子态，不限于初始 Fock 数；交换中的光子产生可以使它增大。
\end{proposition}
\begin{proof}
在全相互作用绘景中把吸收速率差写成加权电子平移
\[
 D_e(t)=\sum_\ell[v_\ell(t)-v_0(t)]
                 |\ell+1\rangle\langle\ell|,
 \qquad R=\ii\hbar(D_ea-D_e^\dagger a^\dagger).
\]
与工作窗相接的各边系数由 $a_W$ 控制，所以
$\|D_eP_W\|,\|D_e^\dagger P_W\|\le a_W$。电子投影与光子算符对易，于是
\begin{align*}
 \frac1\hbar\|RP_W\Psi_0\|
 &\le a_W\bigl(\|aP_W\Psi_0\|
                         +\|a^\dagger P_W\Psi_0\|\bigr)\\
 &\le a_W\bigl(\sqrt{\langle N_{\rm ph}\rangle_{0,t}}
                   +\sqrt{\langle N_{\rm ph}\rangle_{0,t}+1}\bigr)\\
 &\le2a_W\sqrt{\langle N_{\rm ph}\rangle_{0,t}+1}.
\end{align*}
第一步用三角不等式，第二步用 $a^\dagger a=N_{\rm ph}$、$aa^\dagger=N_{\rm ph}+I$ 及 $0\le P_W\le I$。在
\[
 (U-U_0)\Psi_i=-\frac\ii\hbar\int U(t_f,t)R(t)\Psi_0(t)\dd t
\]
中插入 $P_W+(I-P_W)$，再用 $U$ 的幺正性，便得到\eqnrefs{eq:feqo-quantum-ladder-propagation-error}。
\end{proof}

因此相同的逐边失谐在高光子占据下会产生更大的振幅误差；只检查电子反冲频率和作用时间不足以给出所有光态的一致精度。若参考数矩至多为 $\mathcal N$，相对耦合差至多为 $\epsilon_G$，失谐差至多为 $d_\Delta$，则前一项至多为 $2\sqrt{\mathcal N+1}\int|\mathcal G_0(t)|(\epsilon_G+\tau d_\Delta)\dd t$。强脉冲中的累计耦合、玻色增强和工作窗尾部由此都进入同一个可核算的误差界。


在单模、固定纵向动量步长且只保留正能电子通道的有效模型内，对一般光子数混合态 $\rho_{\rm ph}=\sum_n p_n|n\rangle\langle n|$，电子能谱的完整前向模型为
\begin{equation}
\boxed{
 P_\ell
 =\sum_{n=0}^{\infty}p_n\,
 \left|c_\ell^{(n)}(t_f)\right|^2,}
\label{eq:qpinem_finite_recoil_number_forward_dp46}
\end{equation}
其中每一列 $|c_\ell^{(n)}|^2$ 都由同一组物理参数通过\eqnrefs{eq:qpinem_finite_recoil_coupled_lattice_dp46}计算。当逐边差异的传播误差和物理格边界误差均趋于零时，转回全相互作用绘景后的振幅趋于 $\mathcal A_{\ell|n}(g_{\mathrm Q})$，至多差一个与格点无关的 Magnus 相位；联合旋转格还保留共同自由格相位。两种相位都不改变格点概率，因此\eqnrefs{eq:qpinem_finite_recoil_number_forward_dp46}退化为\secref{subsec:single-mode-number-instrument-dp19}使用的条件响应核。

光场趋于经典并不会自动消除电子反冲。在\eqnrefs{eq:qpinem_finite_recoil_coupled_lattice_dp46}中取相干平均场，仍得到含 $\delta_\ell$ 的经典有限反冲格；只有均匀梯传播误差也受控时，才可使用 Bessel 结果。现在在位移模型\eqnrefs{eq:qpinem_S}内定量求这个经典极限。令光场初态为 $|\alpha_{\rm coh}\rangle$，并取
\begin{equation}
 |\alpha_{\rm coh}|\to\infty,
 \qquad |g_{\mathrm Q}|\to0,
 \qquad \beta=\alpha_{\rm coh}g_{\mathrm Q}
 \quad\text{保持有限},
\label{eq:qpinem_coherent_classical_scaling_dp46}
\end{equation}
则 $\hat a=\alpha_{\rm coh}+\delta \hat a$ 使交换生成元分成经典项
$\beta \hat b^\dagger-\beta^*\hat b$ 与涨落项
$g_{\mathrm Q} \hat b^\dagger\delta \hat a-g_{\mathrm Q}^*\hat b\delta \hat a^\dagger$。后者的算符系数按 $|g_{\mathrm Q}|\sim|\alpha_{\rm coh}|^{-1}$ 消失，于是电子约化散射趋向经典 PINEM 算符 $S_{\rm cl}(\beta)$。因此
\begin{equation}
\boxed{
 P_\ell\longrightarrow J_\ell^2(2|\beta|)}
\label{eq:qpinem_to_pinem_probability_limit_dp46}
\end{equation}
这正是固定经典场幅 $\beta$ 时 $|\alpha_{\rm coh}|\to\infty$、$|g_{\mathrm Q}|\to0$ 的强相干态极限。

\begin{proposition}[大相干幅下的经典平移极限]
设光场为相干态，双边电子模型保持 $\beta=\alpha_{\rm coh}g_{\mathrm Q}$ 固定并令 $g_{\mathrm Q}\to0$。位移绘景中经典与涨落生成元严格对易，固定有限涨落 Fock 窗的传播子强收敛到经典平移传播。因此单格入射的约化边带概率趋于\eqnrefs{eq:qpinem_to_pinem_probability_limit_dp46}；此极限不是在微观耦合固定时单独令光强无穷。
\end{proposition}
\begin{proof}
令光场初态为 \(|\alpha_{\rm coh}\rangle\)。用位移算符把大相干幅度移到真空涨落上，
\begin{equation*}
D_{\rm ph}^\dagger(\alpha_{\rm coh})
\hat a
D_{\rm ph}(\alpha_{\rm coh})
=\hat a+\alpha_{\rm coh}.
\end{equation*}
联合散射指数中的生成元因此精确分成
\begin{align*}
X
={}&g_{\mathrm Q}\hat b^\dagger(\hat a+\alpha_{\rm coh})
-g_{\mathrm Q}^*\hat b(\hat a^\dagger+\alpha_{\rm coh}^*)\\
={}&X_{\rm cl}+X_{\rm fl},
\end{align*}
其中
\begin{equation*}
X_{\rm cl}=\beta\hat b^\dagger-\beta^*\hat b,
\qquad
X_{\rm fl}=g_{\mathrm Q}\hat b^\dagger\hat a-g_{\mathrm Q}^*\hat b\hat a^\dagger,
\qquad
\beta\equiv\alpha_{\rm coh}g_{\mathrm Q}.
\end{equation*}
在理想双边电子能量格上
\([\hat b,\hat b^\dagger]=0\)，且电子平移与光子算符彼此对易，因此
\begin{equation*}
[X_{\rm cl},X_{\rm fl}]=0.
\end{equation*}
所以大相干背景与量子涨落不是近似地、而是精确地因子化：
\begin{equation*}
\ee^X
=\ee^{X_{\rm cl}}\ee^{X_{\rm fl}}.
\end{equation*}

现在才取经典场极限
\begin{equation*}
|\alpha_{\rm coh}|\to\infty,
\qquad
|g_{\mathrm Q}|\to0,
\qquad
\beta=\alpha_{\rm coh}g_{\mathrm Q}\ \text{固定}.
\end{equation*}
为说明涨落部分确实消失，而不是只比较两个形式上的系数，取任意固定光子数截断 \(P_N\)。双边平移满足 \(\|\hat b\|=\|\hat b^\dagger\|=1\)，而在有限 Fock 子空间上
\begin{equation*}
\|\hat aP_N\|=\sqrt N,
\qquad
\|\hat a^\dagger P_N\|=\sqrt{N+1}.
\end{equation*}
故
\begin{equation*}
\|X_{\rm fl}P_N\|
\le |g_{\mathrm Q}|\bigl(\sqrt N+\sqrt{N+1}\bigr)
\longrightarrow0.
\end{equation*}
对每个固定的有限涨落扇区，\(\ee^{X_{\rm fl}}\) 因而强收敛到恒等算符；保留下来的电子传播子是
\begin{equation*}
S_{\rm cl}(\beta)
=\exp(\beta\hat b^\dagger-\beta^*\hat b).
\end{equation*}

最后求电子边带振幅。利用正文已经建立的双边平移 Fourier 表示或 Jacobi--Anger 生成函数，
\begin{equation*}
\langle\ell|S_{\rm cl}(\beta)|0\rangle
=J_\ell(2|\beta|)\ee^{\ii\ell\arg\beta}.
\end{equation*}
因此
\begin{equation*}
P_\ell
=\left|\langle\ell|S_{\rm cl}(\beta)|0\rangle\right|^2
=J_\ell^2(2|\beta|),
\end{equation*}
即正文\eqnrefs{eq:qpinem_to_pinem_probability_limit_dp46}。在这一极限中，“经典 PINEM”的数学含义是：相干平均场 \(\beta\) 保持有限，而固定涨落扇区上的量子生成元按 \(|g_{\mathrm Q}|\) 消失。

对涨落生成元 $X_{\rm fl}$，幺正单参数群给 $(e^{X_{\rm fl}}-I)\psi=\int_0^1e^{sX_{\rm fl}}X_{\rm fl}\psi\,\dd s$，故有限数域中 $\|(e^{X_{\rm fl}}-I)\psi\|\le\|X_{\rm fl}\psi\|$，补足由生成元变小到传播子强收敛的一步。
\end{proof}


电子初态的有限纵向带宽则通过输入密度矩阵进入。令
\begin{equation*}
 \rho_\ee^{\rm in}
 =\int\!\dd p\,\dd p'\,
 \rho_e(p,p')|p\rangle\langle p'|,
\end{equation*}
并用总散射算符 $S$ 定义电子约化输出
\begin{equation}
\boxed{
 \rho_\ee^{\rm out}
 =\operatorname{Tr}_{\rm ph}
 \left[S(\rho_\ee^{\rm in}\otimes\rho_{\rm ph})S^\dagger\right].}
\label{eq:finite_envelope_reduced_output_dp47}
\end{equation}
取能量归一化的电子基 $\langle E|E'\rangle=\delta(E-E')$；若还存在未分辨自旋或横向通道，下式的 $p_e(E)$ 对这些指标求和。电子的理想能谱密度为 $p_e(E)=\langle E|\rho_\ee^{\rm out}|E\rangle$，单位为能量的倒数，满足 $\int p_e(E)\dd E=1$。令 $R_E(\varepsilon-E)\ge0$ 表示真实能量 $E$ 被读成 $\varepsilon$ 的条件密度，$\int R_E(x)\dd x=1$。谱仪记录的密度先由卷积给出，再在互不相交的像素能窗 $I_j$ 内积分为无量纲概率：
\begin{equation}
\boxed{
 \begin{aligned}
 p_{\rm det}(\varepsilon)&=\int\dd E\,R_E(\varepsilon-E)p_e(E),\\
 P_j^{\rm meas}&=\int_{I_j}\dd\varepsilon\,p_{\rm det}(\varepsilon).
 \end{aligned}}
\label{eq:pinem_detector_convolution_dp47}
\end{equation}
卷积后的 $p_{\rm det}(E_j)$ 仍是密度，不能直接称为第 $j$ 个像素的概率；窄像素近似为 $P_j^{\rm meas}\simeq p_{\rm det}(E_j)\Delta E_j$。当像素窗覆盖全部读出能量时，交换两次非负积分便得 $\sum_jP_j^{\rm meas}=1$；有限探测范围要另保留漏出概率。后续统计反演应使用\eqnrefs{eq:pinem_detector_convolution_dp47}的像素核，而不是直接把理想 $P_\ell$ 与实验直方图比较。

有限输入能散 $\sigma_E$ 是否会洗掉相邻边带取决于两个不同条件。分辨边带需要
\begin{equation}
 \epsilon_{\rm res}\equiv\frac{\sigma_E}{\hbar\omega}\ll1,
\end{equation}
而保持输入相干的相位条件由输入密度矩阵非对角宽度 $\sigma_{E,{\rm coh}}$ 决定。二者一般不相等；能谱窄并不自动意味着长相干时间。


选择规则 $\ell+m=n$ 是联合 $U(1)$ 对称性的守恒扇区条件，并非某一微扰阶中的偶然零点。

有限反冲模型中的对角失谐来自电子真实色散，非均匀玻色跃迁系数来自光子数矩阵元；两者是独立的结构来源。

\begin{derivation}[从\eqnrefs{eq:multimode_qpinem_continuous_kernel_dp101}到\eqnrefs{eq:multimode_qpinem_lattice_dp47}]

先把第 $\nu$ 个连续电子耦合算符记为
\[
\hat{\mathcal G}_{\nu,I}(t)
=
\sum_{ss'}
\int\dd^3p\,\dd^3p'\,
\mathcal G_{\nu;s's}(\mathbf p',\mathbf p;t)
|\mathbf p',s'\rangle\langle\mathbf p,s|.
\]
在实际窄带工作窗内，假定主导矩阵元集中在一个自旋保持、近固定动量转移通道
\[
(\mathbf p,s_0)
\longrightarrow
(\mathbf p+\hbar\boldsymbol\kappa_\nu,s_0).
\]
定义参考平移算符
\[
\hat b_\nu^\dagger
=
\int\dd^3p\,
|\mathbf p+\hbar\boldsymbol\kappa_\nu,s_0\rangle
\langle\mathbf p,s_0|
\]
以及投影后的共同速率 $\upsilon_\nu(t)$。于是精确耦合算符在该工作窗上分解成
\[
\hat{\mathcal G}_{\nu,I}(t)
=
\ii\upsilon_\nu(t)\hat b_\nu^\dagger
+
\hat R_{\nu,\kappa}(t).
\]
把所有余项收集为
\[
\hat R_\kappa(t)
=
\hbar
\sum_\nu
\left[
\hat R_{\nu,\kappa}(t)\hat a_\nu
+
\hat R_{\nu,\kappa}^\dagger(t)\hat a_\nu^\dagger
\right].
\]
令 $\hat H_{\rm tr}$ 只保留主导平移项，并由它生成 $\hat U_{\rm tr}$。完整传播子 $\hat U$ 与 $\hat U_{\rm tr}$ 的 Duhamel 恒等式为
\[
(\hat U-\hat U_{\rm tr})\Pi_{\rm in}
=
-\frac{\ii}{\hbar}
\int_{t_i}^{t_f}\dd t\,
\hat U(t_f,t)
\hat R_\kappa(t)
\hat U_{\rm tr}(t,t_i)\Pi_{\rm in}.
\]
完整传播子幺正，因此
\[
\|(\hat U-\hat U_{\rm tr})\Pi_{\rm in}\|
\le
\frac1\hbar
\int_{t_i}^{t_f}\dd t\,
\|
\hat R_\kappa(t)
\hat U_{\rm tr}(t,t_i)\Pi_{\rm in}
\|,
\]
这正是正文\eqnrefs{eq:multimode_momentum_step_projection_error_dp101}。只有该量足够小时，连续核才受控退化为
\[
\hat H_{\rm Q}^{(M)}
=
\ii\hbar\sum_\nu
\left[
\upsilon_\nu\hat b_\nu^\dagger\hat a_\nu
-\upsilon_\nu^*\hat b_\nu\hat a_\nu^\dagger
\right].
\]

设整数向量 $\boldsymbol\ell$ 在实际工作窗中唯一对应电子动量
\[
\mathbf p_{\boldsymbol\ell}
=
\mathbf p_0
+\hbar\sum_\nu\ell_\nu\boldsymbol\kappa_\nu,
\]
并定义
\[
|\boldsymbol\ell;\mathbf n\rangle
=
|\mathbf p_{\boldsymbol\ell}\rangle_e
\otimes
\bigotimes_\nu
|n_\nu-\ell_\nu\rangle_\nu .
\]
对第 $\nu$ 个吸收顶角，
\begin{align*}
\hat b_\nu^\dagger\hat a_\nu
|\boldsymbol\ell;\mathbf n\rangle
&=
\hat b_\nu^\dagger
|\mathbf p_{\boldsymbol\ell}\rangle
\otimes
\sqrt{n_\nu-\ell_\nu}\,
|n_\nu-\ell_\nu-1\rangle_\nu
\otimes_{\mu\neq\nu}
|n_\mu-\ell_\mu\rangle_\mu\\
&=
\sqrt{n_\nu-\ell_\nu}\,
|\boldsymbol\ell+\mathbf e_\nu;\mathbf n\rangle .
\end{align*}
同理，
\[
\hat b_\nu\hat a_\nu^\dagger
|\boldsymbol\ell;\mathbf n\rangle
=
\sqrt{n_\nu-\ell_\nu+1}\,
|\boldsymbol\ell-\mathbf e_\nu;\mathbf n\rangle .
\]

在旋去理想多频能量
$E_0+\sum_\nu\ell_\nu\hbar\omega_\nu$
的表象中，真实电子色散留下对角项
\[
\hat H_{\rm rec}
=
\hbar\sum_{\boldsymbol\ell}
\delta_{\boldsymbol\ell}
|\boldsymbol\ell\rangle\langle\boldsymbol\ell|.
\]
展开
\[
|\Psi_{\mathbf n}(t)\rangle
=
\sum_{\boldsymbol\ell}
c_{\boldsymbol\ell}^{(\mathbf n)}(t)
|\boldsymbol\ell;\mathbf n\rangle
\]
并代入
\[
\ii\hbar\partial_t|\Psi_{\mathbf n}\rangle
=
(\hat H_{\rm rec}+\hat H_{\rm Q}^{(M)})
|\Psi_{\mathbf n}\rangle.
\]
左乘 $\langle\boldsymbol\ell;\mathbf n|$ 后，来自相邻格点的两类项分别为
\[
+\ii\hbar
\upsilon_\nu
\sqrt{n_\nu-\ell_\nu+1}\,
c_{\boldsymbol\ell-\mathbf e_\nu}^{(\mathbf n)}
\]
和
\[
-\ii\hbar
\upsilon_\nu^*
\sqrt{n_\nu-\ell_\nu}\,
c_{\boldsymbol\ell+\mathbf e_\nu}^{(\mathbf n)}.
\]
与对角反冲项合并并除以 $\hbar$，得到正文\eqnrefs{eq:multimode_qpinem_lattice_dp47}。

最后说明多指标标签为什么需要一一对应。若存在
$\boldsymbol\ell\neq\boldsymbol\ell'$ 却有
\[
\sum_\nu
(\ell_\nu-\ell_\nu')
\boldsymbol\kappa_\nu=0,
\]
则
$|\mathbf p_{\boldsymbol\ell}\rangle=
|\mathbf p_{\boldsymbol\ell'}\rangle$，
这两个“格点”并非正交电子态。此时必须先把所有到达同一物理动量的路径振幅相加，再定义独立基底；直接把 $\boldsymbol\ell,\boldsymbol\ell'$ 当成两个末态会重复计算 Hilbert 空间维数。
\end{derivation}

\begin{derivation}[从\eqnrefs{eq:multicolor_classical_scattering_dp47}到\eqnrefs{eq:multicolor_pinem_convolution_dp47}]

在多模相干态经典极限中，每个量子模式的算符被分解成平均幅与量子涨落；固定
\[
\beta_\nu=\alpha_\nu\gqmode{\nu}
\]
并令相对涨落修正消失后，电子散射算符为
\[
S_{\rm cl}^{(M)}
=
\exp\!\left[
\sum_\nu
(\beta_\nu\hat b_\nu^\dagger-\beta_\nu^*\hat b_\nu)
\right].
\]
完整连续动量空间中的平移算符彼此对易，
\[
[\hat b_\mu,\hat b_\nu]=
[\hat b_\mu,\hat b_\nu^\dagger]=0,
\]
所以指数可逐模分解：
\[
S_{\rm cl}^{(M)}
=
\prod_\nu
\exp(
\beta_\nu\hat b_\nu^\dagger
-\beta_\nu^*\hat b_\nu).
\]

对单个模式，记
\[
\beta_\nu=|\beta_\nu|\ee^{\ii\varphi_\nu},
\qquad
\hat T_\nu\equiv\hat b_\nu^\dagger.
\]
在理想双边平移空间中 $\hat T_\nu$ 是幺正算符，
\[
\hat T_\nu^{-1}=\hat b_\nu.
\]
因此任意整数幂统一写作
\[
\hat T_\nu^{r}
=
\begin{cases}
(\hat b_\nu^\dagger)^r,&r\ge0,\\
(\hat b_\nu)^{-r},&r<0.
\end{cases}
\]
取 $\hat T_\nu$ 的相位本征态
\[
\hat T_\nu|\theta_\nu\rangle
=
\ee^{-\ii\theta_\nu}|\theta_\nu\rangle,
\]
则单模生成元在该表象中成为普通复数：
\begin{align*}
\beta_\nu\hat T_\nu-\beta_\nu^*\hat T_\nu^{-1}
&\longrightarrow
|\beta_\nu|
\left[
\ee^{\ii(\varphi_\nu-\theta_\nu)}
-\ee^{-\ii(\varphi_\nu-\theta_\nu)}
\right]\\
&=
2\ii|\beta_\nu|
\sin(\varphi_\nu-\theta_\nu).
\end{align*}
Jacobi--Anger 恒等式给
\[
\exp[
2\ii|\beta_\nu|\sin(\varphi_\nu-\theta_\nu)
]
=
\sum_{r_\nu\in\mathbb Z}
J_{r_\nu}(2|\beta_\nu|)
\ee^{\ii r_\nu\varphi_\nu}
\ee^{-\ii r_\nu\theta_\nu}.
\]
把
$\ee^{-\ii r_\nu\theta_\nu}$
重新提升为算符 $\hat T_\nu^{r_\nu}$，得到严格的整数平移展开
\[
\exp(
\beta_\nu\hat b_\nu^\dagger
-\beta_\nu^*\hat b_\nu)
=
\sum_{r_\nu\in\mathbb Z}
J_{r_\nu}(2|\beta_\nu|)
\ee^{\ii r_\nu\varphi_\nu}
\hat T_\nu^{r_\nu}.
\]

对所有模式相乘：
\[
S_{\rm cl}^{(M)}
=
\sum_{\{r_\nu\}}
\left[
\prod_\nu
J_{r_\nu}(2|\beta_\nu|)
\ee^{\ii r_\nu\varphi_\nu}
\right]
\prod_\nu\hat T_\nu^{r_\nu}.
\]
若
\[
\kappa_\nu=m_\nu\kappa_0,
\qquad
m_\nu\in\mathbb Z,
\]
则
\[
\prod_\nu\hat T_\nu^{r_\nu}
\]
把电子总共平移
\[
\Delta p
=
\hbar\kappa_0
\sum_\nu m_\nu r_\nu.
\]
因此所有满足
\[
\ell=\sum_\nu m_\nu r_\nu
\]
的整数路径都到达同一个电子末态 $|\ell\rangle$，其振幅必须相干求和。于是
\[
A_\ell
=
\sum_{\{r_\nu\}}
\delta_{\ell,\sum_\nu m_\nu r_\nu}
\prod_\nu
J_{r_\nu}(2|\beta_\nu|)
\ee^{\ii r_\nu\varphi_\nu},
\]
即正文\eqnrefs{eq:multicolor_pinem_convolution_dp47}。

若不同颜色的动量步长不满足这种整数公度关系，则一般不存在唯一的一维总阶数 $\ell$；此时末态应继续用真实电子动量或多指标通道标记，不能把不同可分辨末态强行卷积到一条 Bessel 能量格。
\end{derivation}



\begin{derivation}[共振单模近场的量子耦合]

一般线性电子--玻色相互作用在相互作用绘景中可写成
\[
\hat V_I(t)
=
\hbar\sum_\lambda
\left[
\hat{\mathcal G}_{\lambda,I}(t)\hat a_\lambda
+
\hat{\mathcal G}_{\lambda,I}^\dagger(t)\hat a_\lambda^\dagger
\right].
\]
这里将光模自由演化的 $\ee^{-\ii\omega_\lambda t}$ 也吸收到 $\hat{\mathcal G}_{\lambda,I}$ 中；$\hat a_\lambda$ 为固定的相互作用绘景模算符，不能再重复附加该相位。$\hat{\mathcal G}_{\lambda,I}$ 仍作用在完整正能电子连续谱上。先把电子自由本征态记为 $|p,s\rangle$，则
\[
\hat{\mathcal G}_{\lambda,I}(t)
=
\sum_{ss'}\int\dd p\,\dd p'\,
\mathcal G_{\lambda;s's}(p',p)\,
\ee^{+\ii\{[E(p')-E(p)]/\hbar-\omega_\lambda\}t}
|p',s'\rangle\langle p,s|.
\]
因此“单模 QPINEM”需要依次做三个投影，而不是一步把完整哈密顿量 改写成升降算符。

第一步是环境谱投影。设目标模式为 $\lambda_0$，其频率为 $\omega_0$；其余模式与目标模式的最小有效失谐记为
\[
\Delta_{\rm spec}
=
\inf_{\lambda\neq\lambda_0}
|\omega_\lambda-\omega_0|.
\]
若作用时间 $T$ 和实际耦合带宽满足
\[
\Delta_{\rm spec}T\gg1,
\qquad
\sup_{\lambda\neq\lambda_0}
\int_{t_i}^{t_f}\dd t\,
\|\hat{\mathcal G}_{\lambda,I}(t)\|
\ll1,
\]
则有限时间 Fourier 核抑制非目标模式，环境和式可约化为 $\lambda_0$。

第二步是电子通道投影。连续谱广义本征态满足 Dirac--$\delta$ 归一化，而离散通道模型需要一组真正 Kronecker 正交的归一态。先在长度 $L_z$ 的周期盒中正则化纵向连续谱，使允许动量间隔为 $2\pi\hbar/L_z$；再选择与主导动量转移 $\hbar\kappa_z$ 相容的一串正交盒归一态。记
\[
|\ell\rangle_V
\equiv
|p_\ell,s_0\rangle_V,
\qquad
p_\ell=p_0+\ell\hbar\kappa_z,
\qquad
{}_V\langle\ell|m\rangle_V=\delta_{\ell m}.
\]
若实际电子是有限宽波包，也可以用彼此近似正交、中心动量为 $p_\ell$ 的归一化通道包代替盒态；最终有效格点只要求相邻通道的重叠误差小于所保留动力学精度。记号上省略盒归一化下标 $V$，而离散化极限保持不变。投影后的电子耦合算符为
\[
\hat{\mathcal G}_{0,I}(t)
=
\sum_\ell
\mathcal G_\ell(t)
\ee^{+\ii\Delta_\ell t}
|\ell+1\rangle\langle\ell|
+
\hat R_{\rm off}(t),
\]
其中
\[
\Delta_\ell
=
\frac{E(p_{\ell+1})-E(p_\ell)}{\hbar}
-\omega_0
\]
是逐边真实失谐，$\hat R_{\rm off}$ 收集自旋翻转、非主导动量转移和其它离散通道。

第三步才是平移对称近似。若显著占据窗口 $|\ell|\le L_{\rm occ}$ 内
\[
\epsilon_{\mathcal G}
\equiv
\max_\ell
\left|
\frac{\mathcal G_\ell-\mathcal G_0}{\mathcal G_0}
\right|
\ll1,
\qquad
\epsilon_\Delta
\equiv
T\max_\ell|\Delta_\ell-\Delta_0|
\ll1,
\]
便可把共同边耦合写成
\[
\mathcal G_\ell(t)\ee^{+\ii\Delta_\ell t}
=
\mathcal G(t)
+
O\!\left(|\mathcal G_0|(\epsilon_{\mathcal G}+\epsilon_\Delta)\right),
\]
并定义双边电子平移算符
\[
\hat b^\dagger
=
\sum_{\ell\in\mathbb Z}
|\ell+1\rangle\langle\ell|,
\qquad
\hat b
=
\sum_{\ell\in\mathbb Z}
|\ell\rangle\langle\ell+1|.
\]
于是
\[
\hat{\mathcal G}_{0,I}(t)
=
\mathcal G(t)\hat b^\dagger
+\delta\hat{\mathcal G}(t).
\]

把这一结果代回单模完整哈密顿量，
\[
\hat V_I(t)
=
\hbar
\left[
\mathcal G(t)\hat b^\dagger\hat a
+
\mathcal G^*(t)\hat b\hat a^\dagger
\right]
+\hat R_{\rm RWA}(t)
+\hat R_{\rm tr}(t),
\]
其中 $\hat R_{\rm tr}$ 收集逐边不均匀和通道投影误差。对一般实线性场，完整相互作用还包含电子上移同时产生光子、电子下移同时湮灭光子的反旋支路；它们在相互作用绘景中带有约 $2\omega_0$ 的快速相位。这些快速支路能否舍弃，不能只由“$\omega_0$ 很大”判断。前章\eqnrefs{eq:rwa_projected_dynamical_error} 已把旋波余项的真实有限时间误差写成 Duhamel 范数；在脉冲包络缓变且反旋相位频率约为 $2\omega_0$ 的常见情形，分部积分才进一步给出\eqnrefs{eq:rwa-finite-time-error-scales-dp68} 那类充分尺度条件。只有该投影传播误差满足
\[
\epsilon_{\rm RWA}^{\rm dyn}\ll1
\]
时，反旋支路才可受控消去；$|\mathcal G|/\omega_0$ 与 $(\omega_0T)^{-1}$ 只是缓变脉冲下对这个更一般误差界的估计。

保留真实矩阵元的记号 $\mathcal G(t)$，并另定义 $\upsilon(t)=-\ii\mathcal G(t)$，即 $\mathcal G(t)=\ii\upsilon(t)$。这样 $g_{\mathrm Q}=\int\upsilon(t)\dd t$ 与经典 $\beta$ 的相位对应是显式的。代入上式，
\[
\boxed{
\hat H_Q(t)
=
\ii\hbar
\left[
\upsilon(t)\hat b^\dagger\hat a
-
\upsilon^*(t)\hat b\hat a^\dagger
\right].
}
\]
这就是正文\eqnrefs{eq:qpinem-single-mode-hamiltonian-dp68}。
\end{derivation}

